\documentclass[11pt]{article}
\usepackage[letterpaper,margin=1in]{geometry}
\usepackage[T1]{fontenc}
\usepackage{lmodern,microtype}
\usepackage{amsmath,amssymb,amsthm,mathtools}
\usepackage{booktabs,longtable,array}
\usepackage{xcolor,xurl}
\usepackage[colorlinks=true,linkcolor=blue!45!black,citecolor=blue!45!black,urlcolor=blue!45!black]{hyperref}
\setlength{\emergencystretch}{2em}
\newtheorem{theorem}{Theorem}[section]
\newtheorem{lemma}[theorem]{Lemma}
\newtheorem{proposition}[theorem]{Proposition}
\newtheorem{corollary}[theorem]{Corollary}
\theoremstyle{definition}
\newtheorem{definition}[theorem]{Definition}
\theoremstyle{remark}
\newtheorem{remark}[theorem]{Remark}
\title{\textbf{Subexponential Extended Frege refutations of unsatisfiable 3-CNFs}\\[0.4em]
\large relative to an integral grid-contraction construction}
\author{Charlie Cruz}
\date{19 November 2026, revised 25 November 2026}
\hypersetup{pdftitle={Subexponential Extended Frege refutations of unsatisfiable 3-CNFs},pdfauthor={Charlie Cruz}}
\begin{document}
\maketitle

\begin{abstract}
We take as an input the supplied integral grid-contraction
construction, including its local identities, formula construction,
and index bounds. Relative to this input, every unsatisfiable 3-CNF
with $n$ variables and $m$ clauses has an Extended Frege refutation
of size $2^{O((n+m)/\sqrt{\log(n+m+2)})}$. Proof-efficient
sparsification then gives $2^{o(n)}$-size Extended Frege refutations
of all unsatisfiable normalized $n$-variable 3-CNFs. The argument does
not assume that Extended Frege simulates general IPS or proves the
correctness of general PIT.

The proof attaches an independent pair of weights to every edge
endpoint of the constraint network after its initial subdivision.
Pinning these pairs to a Boolean assignment makes each initial clique
table a scalar times a product of functions of individual incident
labels, and every fixed grade-and-split merge of the contraction
preserves this form. This yields a short arithmetic derivation of the
pinning identity, whose size and exact integer bit widths are both
$2^{O((n+m)/\sqrt{\log(n+m+2)})}$. A conditional-expectation argument
with shared Boolean selectors converts the pinning identity and a zero
count into an Extended Frege refutation. The same ingredients give
CLIP plus Extended Frege proofs of exact model counts within the same
size bound.

The corresponding statement for \emph{ordinary Frege} remains
unproved. The additional obstacle is removing the sequential extension
definitions without losing the subexponential size bound. Two
appendices prove obstructions to particular replacement strategies:
normalizing coefficients in a single variable order, even an
adaptively chosen one, and certifying positivity by explicit real sums
of polynomial squares.
\end{abstract}

\tableofcontents
\bigskip

\section{Introduction}
\label{sec:introduction}

Szeider~\cite{szeider} (version~2, \S7, Open Problem~1, printed
p.~10) asks whether every unsatisfiable 3-CNF has a $2^{o(n)}$-size
Extended Frege (EF) refutation without the hypothesis $(S_p)$ that EF
polynomially simulates IPS~\cite{gp} over an odd-prime field. The paragraph that
follows it (printed p.~11) explains why a soundness invariant asserting
that every signed intermediate is a partial model count fails. The
invariant used here, edge factorization
(Section~\ref{sec:direct-ef}), permits arbitrary signed factors. We
answer this version of the question relative to the supplied
contraction construction~\cite{cruz-ntime}. The stronger
ordinary-Frege question is not answered.

\subsection{Main results}

The integral local identities, formula construction, and index bounds
of the supplied manuscript~\cite{cruz-ntime} are imported as proved
inputs (Section~\ref{sec:imported}). Relative to them we prove the
following.
\begin{enumerate}
\item Every unsatisfiable 3-CNF with $n$ variables and $m$ clauses
has an EF refutation of size $2^{O((n+m)/\sqrt{\log(n+m+2)})}$; in
particular, every fixed-density family has refutations of size
$2^{O(n/\sqrt{\log n})}$ (Theorem~\ref{thm:direct-subexp-ef}).
\item Every unsatisfiable normalized $n$-variable 3-CNF has an EF
refutation of size $2^{o(n)}$ (Theorem~\ref{thm:all-3cnf-ef}). This
follows from the first result and the sparsification tree of
Impagliazzo, Paturi, and Zane~\cite{ipz}. Corollary~\ref{cor:fixed-width}
extends the variable-only bound to every fixed clause width.
Theorem~\ref{thm:quant-width} gives an explicit rate and also covers
widths $k$ with $k\log(k+2)=o(\log\log n)$.
\item Every 3-CNF with $n$ variables and $m$ clauses has a CLIP plus
EF proof of its exact model count, in the sense of Chede, Chew, and
Shukla~\cite{ccs}, of size $2^{O((n+m)/\sqrt{\log(n+m+2)})}$
(Theorem~\ref{thm:clip-counting}).
\item The refutations can be output within the same time bound
(Corollary~\ref{cor:ef-proof-output}), and the same upper bounds hold
for Extended Resolution (Section~\ref{sec:ef-boundaries}).
\end{enumerate}
The new derivation is a mathematical proof with explicit proof-size
accounting. It has not been translated into Lean or into a generated
EF proof file. This does not change which hypotheses the theorems use.

\subsection{What is imported}
\label{sec:imported}

From the supplied manuscript \texttt{separating-ntime-from-space.tex}
\cite{cruz-ntime} we use \texttt{lem:codes}, \texttt{lem:top},
\texttt{lem:atone}, \texttt{cor:cuts}, \texttt{lem:transport},
\texttt{lem:merge}, \texttt{lem:paths}, and the explicit formulas in
\texttt{sec:merge} and \texttt{sec:formula}. (These are labels in that
source file.) The input includes the fact that the orientation is
selected from the entire grade list before summing over splits.

From the literature we use elementary binary-arithmetic EF
proofs~\cite{kls} (Section~\ref{sec:integer-ef-interface}) and the
IPZ sparsification tree~\cite{ipz} (Section~\ref{sec:ef-sparsification}).
Exact locators are given in those sections.

\subsection{What remains open}

The corresponding statement for ordinary Frege, with total written
proof size, is not proved. In the negative-descent argument each
selector bit is an extension variable defined from the earlier
selectors, and a bound on EF proof size does not bound the formulas
obtained by expanding these definitions.
Lemma~\ref{lem:extension-depth-transfer} shows that extension
dependency depth $o(n)$ would suffice; the present proof does not
establish such a bound (Section~\ref{sec:ef-boundaries}).
Appendices~\ref{sec:sos-length-obstruction}
and~\ref{sec:all-orders-coefficient-barrier} prove obstructions to two
particular replacement strategies. No Boolean circuit lower bound and
no separation between Frege and EF is claimed.

\subsection{Overview of the proof}

Section~\ref{sec:integer-ef-interface} fixes the arithmetic interface:
an explicitly listed integer derivation whose expressions have bounded
full-result bit width translates into EF with polynomial overhead
(Lemmas~\ref{lem:integer-derivation-ef}
and~\ref{lem:formula-integer-width}).

Section~\ref{sec:direct-ef} contains the main argument. We start from
the Boolean constraint network of the 3-CNF, apply the supplied
initial-clique construction, and only then attach a weight pair
$(w_{p,0},w_{p,1})$ to every port $p$. Let $W(\mathbf w)$ be the
compiled contraction formula. Substituting the indicator weights
$\operatorname{pin}(a)$ of a Boolean port assignment $a$ makes every
initial clique table \emph{edge-factorized}: a scalar times a product
of one factor per incident pair label
(Lemma~\ref{lem:initial-port-pinning}). Edge factorization survives
each merge with fixed grades and splits
(Lemma~\ref{lem:rankone-closure}), and after the split sum the result
no longer depends on the orientation chosen for those grades
(Lemma~\ref{lem:rankone-fixedgrade}). Induction on the recursion gives
the pinning identity $W(\operatorname{pin}(a))=f_{\rm net}(a)$
(Corollary~\ref{cor:port-pinning-identity}). The negative-descent
lemma (Lemma~\ref{lem:ef-negative-descent}) then converts the pinning
identity into the witness inequality $W(\mathbf1)\ge f(a)$: if it
failed, a sequence of extension-variable selectors would find a
Boolean $b$ at which the pinned value is negative, which the pinning
identity forbids. Together with the numerical fact $W(\mathbf1)=0$
for an unsatisfiable input, this gives the refutation.

Section~\ref{sec:rankone-size-audit} bounds the size and bit widths of
the pinning derivation by $2^{O(s/\sqrt{\log s})}$.
Section~\ref{sec:ef-sparsification} removes the dependence on the
number of clauses by sparsification, and
Section~\ref{sec:clip-counting} derives the CLIP counting proofs.
Section~\ref{sec:ef-boundaries} records consequences and the remaining
ordinary-Frege problem. Appendix~\ref{sec:sos-length-obstruction}
shows that explicit sums of squares cannot replace the signed formula
in general, and Appendix~\ref{sec:all-orders-coefficient-barrier}
shows that coefficient normalization fails in every variable order.
Appendix~\ref{sec:verification} records the checks that were made.

\subsection{Conventions}

All logarithms are base two. Small parameter values are covered by
increasing absolute constants. A 3-CNF has clauses of width at most
three; clauses are treated as sets and duplicate clauses are removed.
Consequently a formula on $n$ variables has $O(n^3)$ distinct
non-tautological clauses. If a literal presentation contains
arbitrarily many repetitions and the proof must write its entire
negation, an additional polynomial term in that presentation length
must be included. A refutation can equivalently be given from the
clauses or as a proof of the negated conjunction; the deduction
conversion has polynomial overhead. Fresh extension variables are
renamed when separate proofs are combined.

\subsection{Novelty and prior work}

We checked the supplied Szeider note and its November 2026 announcement
on the ETH discussion site, the current ECCC 2026 list, and relevant
STOC 2027 submissions. These checks located no unconditional EF upper
bound of this form. This was a bounded literature search; it does not
rule out independent work. The arithmetic library~\cite{kls},
the sparsification method~\cite{ipz}, and the Gram-matrix lower-bound
method~\cite{clr} are credited to existing work.

\section{Integer arithmetic in Extended Frege}
\label{sec:integer-ef-interface}

This section states the single arithmetic input we import and derives
the form in which the later sections use it:
Lemma~\ref{lem:integer-derivation-ef} translates explicit integer
derivations into EF, and Lemma~\ref{lem:formula-integer-width} bounds
the bit widths that the translation needs.

\paragraph{Proof-size conventions.}
We use Extended Frege in its usual finite propositional basis, with
fresh extension variables naming Boolean gates. Size means total written
symbol length. Circuit-style gate and inference counts differ from this
measure by a polynomial factor, because gate identifiers have
logarithmic length and each gate definition has constant fan-in.

\paragraph{The arithmetic library.}
The only arithmetic proof-complexity input is the following publicly
stated proof library. Krapivin, Li, and Shi~\cite{kls} (STOC 2027
submission 1284, Appendix~A.1, printed pp.~43--44) specify full-result
signed binary addition, subtraction, negation, and schoolbook
multiplication. Their Lemma~C.1 (printed pp.~63--65) gives
polynomial-width EF derivations of the addition identities~(62),
multiplication identities~(63), distributivity~(64)--(67), and order
rules~(68)--(76). We take scale zero in that lemma, so that its words
represent integers. The word-width parameter in its statement is an
arbitrary bound on the operand and intermediate word widths; the name
``short-word'' does not restrict that quantifier.

This input proves finitely many elementary arithmetic schemata. It is
not a reflection principle for arbitrary arithmetic identities or for
the counting algorithm. The cited lemma obtains its polynomial bounds
from Cook's propositional translation~\cite{cook75}, stated precisely as Theorem~2.3
of Arteche, Atserias, de Rezende, and Khaniki~\cite{aadk} (June
2025 version, printed p.~18, PDF p.~20): every fixed universal
quantifier-free theorem of $\mathrm{PV}_1$ has polynomial-time
constructible EF proofs of its propositional translations. Krapivin,
Li, and Shi explicitly identify the binary circuits and the elementary
schemata to which this fact is applied. Their stronger quasilinear
theorem is not needed here.

\paragraph{Source check.}
We obtained the complete 148-page STOC submission and inspected its
Appendix~A.1, Lemma~C.1, Theorem~A.2, and Appendix~J.3. Lemma~C.1
states the polynomial arithmetic library and invokes the known
bounded-arithmetic translation; it does not print every bit-level
derivation. We import this existing library as a theorem. The counting
argument does not assume that an arbitrary polynomial-time calculation
has a short correctness proof.

\begin{lemma}[Explicit integer derivations translate to EF]
\label{lem:integer-derivation-ef}
Suppose an explicitly listed arithmetic derivation has total description
length at most $T$. Its terms are given by acyclic shared definitions
using integer constants, Boolean input bits, addition, subtraction, and
multiplication. Every step is one of the following:
\begin{enumerate}
\item an instance of a ring identity or an elementary order rule from
the finite library described above;
\item equality substitution, congruence, transitivity, or propositional
case analysis;
\item a ground integer identity supplied with a complete arithmetic
calculation;
\item an explicitly displayed finite sum or product manipulation
expanded into the preceding kinds of steps.
\end{enumerate}
Assume every arithmetic expression used, including expressions in
ground calculations and local rewrite instances, has a full-result
binary implementation of width at most $B$. Then the derivation has
an EF translation of size $(T+B+1)^{O(1)}$. The exponent is an absolute
constant for the fixed arithmetic library. No polynomial identity
testing axiom is used.
\end{lemma}
\begin{proof}
For every input bit and constant, introduce the needed signed binary
word, and for every arithmetic operation introduce the Boolean gates
of its binary implementation in topological order. An extension
definition uses only earlier variables. An addition or comparison
uses $O(B)$ gates, and the schoolbook multiplier uses $O(B^2)$ gates.
Thus all computation definitions together have size polynomial in
$T+B$. Equality of two words is the conjunction of their bitwise
equalities after sign extension to a common width.

For a ring or order step, instantiate the corresponding fixed library
derivation on the operand words. The operand bits are existing
variables, so this operation does not substitute exponentially large
Boolean formulas for them. The library cost is polynomial in $B$.
It includes any auxiliary gate definitions and their equality checks.
Rename its new variables before inserting the instance, so that all
extension definitions remain fresh and acyclic. There are at most $T$
instances, and their total cost is polynomial in $T+B$.

Congruence can also be implemented directly. If two input tuples to a
Boolean gate agree, a fixed tautology for that gate derives agreement
of its outputs. Traverse the gates of the relevant arithmetic circuit
in topological order and concatenate these constant-size instances.
This proves that equal operand words give equal output words using a
number of steps polynomial in $B$. Transitivity and propositional
case analysis use fixed tautology templates. A conjunction of bitwise
equalities is formed or projected using one template per bit.

In a ground calculation, every input and every gate output is a
specified bit. For each Boolean gate, its fixed truth table yields a
constant-size derivation of the stated output bit from its input bits
and gate definition. Topological traversal verifies the complete
calculation. Consequently a true, explicitly computed ground integer
identity costs polynomially many EF symbols in its calculation length
and bit width; it is not admitted as an additional axiom.

An explicitly listed reordering of a finite sum or product can be
carried out by adjacent transpositions and reassociation. If there
are $r\le T$ entries, at most $r^2$ transpositions suffice. Each
transposition uses a constant number of associativity and
commutativity instances. Distribution over an explicitly listed index
set uses one distributivity instance per written term, followed by
such reordering steps. The hypothesis charges all intermediate
expansions to $T$, so these operations also have polynomial total
cost. This explains why a succinctly specified sum with $2^n$ terms
cannot be inserted at this stage unless its expansion is actually
within the stated budget.

Finally, full-result widths are assigned to the concrete expressions,
not guessed from their semantic values. For signed inputs of widths
$r,s$, addition uses $\max(r,s)+1$ bits, negation uses $r+1$
bits, and multiplication uses $r+s$ bits. These lengths represent
the exact integer results for all inputs of the declared widths.
Comparisons sign-extend without changing values. The width hypothesis
therefore rules out modular wraparound throughout the derivation.
Combining the preceding polynomial bounds proves the lemma.
\end{proof}

\begin{lemma}[Formula size controls the necessary bit widths]
\label{lem:formula-integer-width}
An arithmetic formula with at most $t$ leaves, each of signed width
at most $b_0$, and at most $t'$ unary or binary operation gates has
full-result implementations of width $O((t+t')(b_0+1))$ at every
gate. The same bound holds when its leaves are arbitrary Boolean
functions of external inputs.
\end{lemma}
\begin{proof}
At a leaf, reserve $b_0+1$ bits. At an addition or subtraction,
the new width is at most the maximum of the child widths plus two;
this is at most their sum plus two. At a multiplication it is at
most their sum, and at a negation it is the child width plus one.
Induction on the formula therefore bounds its output width by the
sum of its leaf widths plus twice its operation count. The same
induction applies to every subformula. The argument uses only the
range of each leaf, so a Boolean function in place of a Boolean
leaf has the same one-bit range.
\end{proof}

\paragraph{Example: one order rule at the bit level.}
The negative-descent step in the proof of
Lemma~\ref{lem:ef-negative-descent} needs the implication
$u+v<0\ \wedge\ u\ge0\Rightarrow v<0$. If $v\ge0$ as well,
the two inputs have zero sign bits. First align their old signed
widths. In that old highest column, both input bits are zero, and
the incoming carry is at most one, so the outgoing carry is zero.
Now sign-extend by one extra bit. The new high column receives
two zero input bits and the zero carry just established, so its
output sign bit is zero. This contradicts the sign bit asserting a
negative sum. This verification uses a fixed gate argument per
column and also explains why omitting the extra bit would be wrong.
For example, four-bit two's-complement $7+1$ must be evaluated in
five bits as $01000$, not truncated to the negative four-bit word
$1000$.



\section{Port weights, pinning, and negative descent}
\label{sec:direct-ef}

This section proves the main theorem, Theorem~\ref{thm:direct-subexp-ef},
for \emph{Extended Frege}. It does not prove the corresponding
statement for ordinary Frege. The input is the integral contraction
compiler of the supplied manuscript~\cite{cruz-ntime}, including its
local crossing identities and its index bounds. We prove that the
compiler is sound on a restricted class of inputs, the edge-factorized
cliques, which is closed under its recursion. We then use this
soundness proof to establish the witness inequality
(Lemma~\ref{lem:ef-negative-descent}) needed for refutations.

All arithmetic identities below are identities over the integers.
When translated into propositional proofs, integers are written in
binary, with enough bits to avoid overflow. The arithmetic translation
is Lemma~\ref{lem:integer-derivation-ef}, and the bound on bit lengths
is Lemma~\ref{lem:rankone-derivation-size} in
Section~\ref{sec:rankone-size-audit}.

\subsection{Weights at every port after subdivision}

Start with the bounded-degree Boolean network for a 3-CNF. Its vertices
are clause predicates and equality predicates. A variable that occurs
several times is connected through a tree of equality predicates of
degree at most three. Hence the network has $O(n+m)$ edges and vertices,
and the assignment of an original variable determines every edge in
its equality tree. For a variable occurring once, use a degree-one
predicate identically one at the other endpoint of its occurrence
edge. For a specified variable with no occurrences, include a separate
edge between two degree-one predicates identically one. It contributes
the two choices of that variable. There are no auxiliary choices. A satisfying
assignment of the 3-CNF therefore has exactly one satisfying extension
to the network.

Apply the initial-clique construction of the supplied manuscript,
including its subdivision of every network edge by an equality vertex.
Only \emph{after this subdivision}, put an independent pair of weights
\[
 (w_{p,0},w_{p,1})
\]
at every port $p$ of the subdivided network. If the original network has
$e_{\rm net}$ edges, there are $N=4e_{\rm net}$ such ports. Every port has one owner, namely
its incident vertex. Replace the table of a vertex $v$ by
\[
 f_v(z)\prod_{p\in\operatorname{ports}(v)}w_{p,z_p}.
 \tag{\thesection.1}\label{eq:port-weight-table}
\]
No weights are attached to padding requests in the label encoding.
At all-one weights the value is the original model count. Write
$W(\mathbf w)$ for the contraction formula that the supplied compiler
produces from these weighted tables.

For $a\in\{0,1\}^N$, write
\[
 \operatorname{pin}(a)_{p,b}=[a_p=b],
 \qquad
 M_a(\mathbf w)=\prod_{p=1}^N w_{p,a_p}.
\]
The Boolean predicate on port assignments is
\[
 f_{\rm net}(a)=
 \prod_v f_v(a|_{\operatorname{ports}(v)})
 \prod_{e=uv}[a_{u,e}=a_{v,e}].
 \tag{\thesection.2}\label{eq:port-predicate}
\]
The products in this expression are conjunctions of bounded-arity
Boolean predicates. In particular, $f_{\rm net}$ has a Boolean formula
of size polynomial in $n+m$.

\begin{lemma}[Smooth scopes for the weighted compiler]
\label{lem:port-smoothness}
The weighted contraction formula $W(\mathbf w)$ has a scope at every
gate with the following properties. Every nonzero monomial at a gate
contains exactly one weight from each pair in that gate's scope and
no weights from the other pairs. The children of a sum have the same
scope. The children of a product have disjoint scopes whose union is
the parent's scope. The output scope is all $N$ port pairs.
The formula has size $2^{O(s/\sqrt{\log s})}$ for $s=O(n+m)$.
\end{lemma}
\begin{proof}
At an original table entry, the scope is the set of ports owned by
that table. Equation~\eqref{eq:port-weight-table} contains exactly one
weight from each such pair. Grouping vertices takes products over
disjoint vertex sets, hence disjoint port sets. A merged table entry is
a sum of products containing one old entry from every node in its row
or column. Those old entries have disjoint scopes. Equality tests,
weight tests, routing coefficients, and packing tests contain no port
weights. Thus every product summand has exactly the union of the old
scopes, independent of its summation indices. Give zero terms the
declared scope of the expression that they replace. This proves the
assertion inductively through all merges. At the last clique, each
product contains the entries of its two nodes and therefore all
original scopes. The top-level sums preserve this scope.

Replacing an original Boolean table value by the monomial in
\eqref{eq:port-weight-table} costs a polynomial factor in $s$ per
occurrence. The supplied compiler has size
$2^{O(s/\sqrt{\log s})}$, so this factor is absorbed in the bound.
\end{proof}

\begin{lemma}[Pinning the initial clique]
\label{lem:initial-port-pinning}
After substituting $\mathbf w=\operatorname{pin}(a)$, every table of
the initial clique has the form
\[
 T_v((\theta_{vu})_{u\ne v};a)
   =c_v(a)\prod_{u\ne v}b_{vu}(\theta_{vu};a).
 \tag{\thesection.3}\label{eq:initial-rankone}
\]
Here $c_v(a)$ is the product of the predicates of the subdivided-network
vertices in group $v$, and $b_{vu}$ is the indicator of the unique
valid encoded label prescribed by the pins at that endpoint of the
pair. The identities have propositional proofs of size polynomial in
the number of explicitly represented initial table entries.
\end{lemma}
\begin{proof}
For a bounded-arity vertex, the equality
\[
 f_v(z)\prod_{p\in\operatorname{ports}(v)}[z_p=a_p]
 =f_v(a|_{\operatorname{ports}(v)})
   \prod_{p\in\operatorname{ports}(v)}[z_p=a_p]
 \tag{\thesection.4}\label{eq:vertex-pin}
\]
holds for every setting of its bits. If one equality test fails, both
sides vanish. If all tests hold, the two arguments of $f_v$ are the
same. Its arity is bounded by a constant, so a truth-table proof of
\eqref{eq:vertex-pin} has constant size before substitution.

Multiply these identities over the vertices in an initial clique
group. There is no subdivided-network edge internal to a group, by
the initial-clique construction. Regroup the equality tests by the
neighboring clique group. For a fixed neighbor, their product is the
indicator of the encoded bit string prescribed by the pins. For a
label outside the encoding's image, the original table is zero and
this indicator is defined to be zero. This proves
\eqref{eq:initial-rankone}. One bounded-arity identity is used per
original table factor, followed by products and reassociations.
Consequently the proof length is polynomial in the explicitly written
table descriptions.

Pinning only the pre-subdivision network would not prove this lemma:
an unpinned subdivision equality vertex can correlate distinct clique
edges. This is why the weights were attached after subdivision.
\end{proof}

\subsection{Edge-factorized cliques}

We now isolate a class of cliques that the signed recursion maps to
itself and on which the compiler's output can be computed exactly.
Call a clique \emph{edge-factorized} when its tables have the form
\[
 T_v=c_v\prod_{u\ne v}b_{vu}(\theta_{vu}).
 \tag{\thesection.5}\label{eq:edge-factorized}
\]
The factors may be signed. Absorb the endpoint-view reversal into the
notation $b_{vu}$, so that a pair has one common label domain $D_{uv}$.
All summation variables below denote these canonical pair labels.
Whenever a routing function such as $N^{\uparrow}$,
$N^{\downarrow}$ or $\Pi_k$ is written on such variables, it
is understood to be precomposed, at each corner, with the
canonical-label-to-local-view map: the earlier endpoint keeps
the two multisets in their order and the later endpoint swaps
them. This convention applies to both row and column cuts.
Consequently the crossing in canonical coordinates is the literal
equality product $[\xi=\nu][\eta=\mu]$. For example, in the
row cut the displayed lower factor
$N^{\downarrow}(\mu,\nu;m)$ means the manuscript's local-view
function evaluated at $(\operatorname{rev}\mu,
\operatorname{rev}\nu;m)$.
Put
\[
 B_{uv}=\sum_{\theta\in D_{uv}}b_{uv}(\theta)b_{vu}(\theta),
 \qquad
 V(\mathcal C)=\left(\prod_vc_v\right)\left(\prod_{\{u,v\}}B_{uv}\right).
 \tag{\thesection.6}\label{eq:factorized-value}
\]
The value of an edge-factorized clique is $V(\mathcal C)$: each pair
label occurs in only its two endpoint factors, and the pair-label
domains are independent, so distributing the finite sums gives
\eqref{eq:factorized-value}.

\begin{lemma}[Edge factorization survives a fixed merge]
\label{lem:rankone-closure}
Fix a grade list $\mathbf k$, choose the compiler's orientation for
that list, and fix a split list $\mathbf c$. If the parent is
edge-factorized, then the child clique is edge-factorized. Its factors
have explicit formulas using sums over the labels of one old pair or
the two corner views of one rectangle.
\end{lemma}
\begin{proof}
First consider the row orientation. For a row $A$, define
\[
 c'_A=
 \prod_{v\text{ in row }A}c_v
 \prod_{\substack{\{u,v\}\text{ horizontal}\\u,v\text{ in row }A}}B_{uv}.
 \tag{\thesection.7}\label{eq:rankone-newscalar}
\]
For a rectangle $R$ with corners $(x,y;u,v)$, upper row $A$ and lower
row $A'$, define
\begin{align*}
 G_{R,A}(m)
 &= [\operatorname{wt}(m)\le H_R+k_R/2]
    \sum_{\xi,\eta}
       b_{xv}(\xi)b_{yu}(\eta)
       [z^{k_R-c_R}]N^{\uparrow}(\xi,\eta;m),\\
 G_{R,A'}(m)
 &= [\operatorname{wt}(m)\le H_R+k_R/2]
    \sum_{\mu,\nu}
       b_{uy}(\mu)b_{vx}(\nu)
       [z^{c_R}]N^{\downarrow}(\mu,\nu;m).
 \tag{\thesection.8}\label{eq:rankone-newrectangle}
\end{align*}
Each summation uses the old pair-label domains, with the manuscript's
endpoint-view convention. On an unpacked tuple $\chi_{A,A'}$, let
$b'_{A,A'}(\chi_{A,A'})$ be the product of the endpoint factor of every
vertical label stream and the factor $G_{R,A}(m_R)$ of every rectangle
stream on these two rows. Define the factor on a packed label by
unpacking. Set it to zero when packing is invalid or any individual
stream lies outside its prescribed domain.

In the merged-table sum, an old copied label belongs to exactly one
of three disjoint groups. A horizontal pair contributes its two
copies, its two endpoint factors, and an equality test; summing the
copies gives $B_{uv}$. A vertical pair contributes one copy at this
row, its endpoint factor, and a test against the stream label; its sum
gives that endpoint factor at the stream label. A rectangle contributes
its two corner views at this row, their endpoint factors, and the
coefficient of $N^{\uparrow}$ or $N^{\downarrow}$; its sum is the
corresponding expression in~\eqref{eq:rankone-newrectangle}. There are
no other factors involving these copies. Distribute the finite sums
over these disjoint groups. The result is
\[
 T'_A=c'_A\prod_{A'\ne A}b'_{A,A'}.
\]
The packing and stream-domain tests involve one neighbor label each,
so they preserve this factorization. No division by $c'_A$ occurs;
the argument includes zero scalar factors.

For a column merge, exchange rows and columns. Use message budget
$H_R-k_R/2$ and coefficient exponents $-k_R-c_R,c_R$, as in the
compiler's column factorization. The same partition of copied labels
into disjoint groups proves the asserted factorization.
\end{proof}

For the same rectangle, define its weighted grade contraction by
\[
 P_R(k)=\sum_{\xi,\eta,\mu,\nu}
 b_{xv}(\xi)b_{yu}(\eta)b_{uy}(\mu)b_{vx}(\nu)
 \Pi_k(\xi,\eta,\mu,\nu).
 \tag{\thesection.9}\label{eq:weighted-grade}
\]

\begin{lemma}[The orientation disappears after the split sum]
\label{lem:rankone-fixedgrade}
For a fixed grade list $\mathbf k$, sum the values
$V(\mathcal C^{o(\mathbf k)}_{\mathbf k,\mathbf c})$ of the
factorized child cliques over all split lists. The result is
\[
 \left(\prod_vc_v\right)
 \left(\prod_{\text{straight pairs }\{u,v\}}B_{uv}\right)
 \left(\prod_R P_R(k_R)\right),
 \tag{\thesection.10}\label{eq:rankone-fixedgrade}
\]
whether $o(\mathbf k)$ is the row or the column orientation.
\end{lemma}
\begin{proof}
For the row orientation, expand one new edge inner product using its
unpacked tuple. The packing map is a bijection between valid packed
labels in the domain and the allowed tuples. Invalid labels contribute
zero. The tuple streams are independent, so its inner product is the
product of the inner products of the vertical streams and the scalar
contractions of its rectangle streams. Together with the new scalar
factors in~\eqref{eq:rankone-newscalar}, every straight old pair occurs
once as $B_{uv}$.

For a rectangle stream, sum over its split and its allowed message.
The row-cut identity for $\Pi_{k_R}$ says that the sum of the two
coefficient factors, summed over the split and the message, is
$\Pi_{k_R}$ at the four old views. The heavier
messages removed by the weight tests have zero product by the
compiler's row-cut bound. Multiplying by the four old endpoint factors
and distributing their finite sums gives $P_R(k_R)$ from
\eqref{eq:weighted-grade}. Different rectangles have independent split
indices, so their sums factor into the product in
\eqref{eq:rankone-fixedgrade}.

The column-cut identity gives precisely the same $\Pi_{k_R}$, and
therefore the same scalar expression, for the column orientation.
The orientation is fixed by the grades before any split is summed.
This dependence is necessary: the proof would not justify choosing a
different orientation separately for each split term.
\end{proof}

\begin{lemma}[Correctness on edge-factorized tables]
\label{lem:rankone-correctness}
The scalar output of the signed contraction compiler on an
edge-factorized clique equals $V(\mathcal C)$ from
\eqref{eq:factorized-value}. It has an arithmetic derivation by
induction on the number of remaining merges, using the local
crossing identities, finite distributivity, and the identities in
Lemmas~\ref{lem:rankone-closure}--\ref{lem:rankone-fixedgrade}.
\end{lemma}
\begin{proof}
With two nodes, the compiler sums the product of the two entries over
their one shared label. Substituting~\eqref{eq:edge-factorized} gives
$c_uc_vB_{uv}$, which is~\eqref{eq:factorized-value}.

At a larger clique, fix its grades and splits. Lemma
\ref{lem:rankone-closure} gives factorized child tables. Apply the
inductive identity to that child output, then sum over its splits.
Lemma~\ref{lem:rankone-fixedgrade} gives
\eqref{eq:rankone-fixedgrade}. The resulting expression is independent
of which orientation was selected for those grades.

Sum over all grade lists. Their domains are independent across
rectangles, hence finite distributivity gives
\[
 \left(\prod_vc_v\right)
 \left(\prod_{\text{straight pairs}}B_{uv}\right)
 \prod_R\left(\sum_kP_R(k)\right).
\]
The crossing identity
$\sum_k\Pi_k(\xi,\eta,\mu,\nu)=[\xi=\nu][\eta=\mu]$
implies, upon substitution into~\eqref{eq:weighted-grade},
\[
 \sum_kP_R(k)=B_{xv}B_{yu}.
\]
Every non-straight pair is a diagonal of exactly one rectangle.
Thus the last displayed expression is $V(\mathcal C)$.

All these equalities hold for arbitrary signed endpoint factors.
The proof never identifies a signed intermediate entry with a
nonnegative model count.
\end{proof}

\begin{corollary}[The pinning identity]
\label{cor:port-pinning-identity}
The weighted compiled formula satisfies
\[
 W(\operatorname{pin}(a))=f_{\rm net}(a).
 \tag{\thesection.11}\label{eq:compiled-pinning}
\]
Its proof is the initial pinning identities followed by the
factorized-correctness derivation.
\end{corollary}
\begin{proof}
Use Lemma~\ref{lem:initial-port-pinning} and then
Lemma~\ref{lem:rankone-correctness}. An initial pair factor is an
indicator of a prescribed encoded label. Its endpoint inner product
is one exactly when the two prescribed labels agree, and zero
otherwise. The injectivity of the string-to-label encoding makes this
equivalent to agreement of the corresponding port bits. This local
simplification has a proof over that single pair's label domain:
for agreeing prescribed strings exactly one encoded label has both
indicators one, and for disagreeing strings no label does. One-hotness
of the pin indicators follows by expanding
$\prod_i(a_i+(1-a_i))=1$; a disagreement makes the corresponding
factor zero. This uses only the single-edge domain bounded in
Lemma~\ref{lem:rankone-local-domains} of
Section~\ref{sec:rankone-size-audit}. That lemma uses only the
compiler's parameters and bounds (the potential bound, label-code
lengths, and grade and split ranges) and no result of the present
section, so this forward reference is not circular.
The scalar factors are exactly the vertex-predicate factors in
\eqref{eq:port-predicate}. Their product is $f_{\rm net}(a)$.
\end{proof}

\subsection{From pinning to the witness inequality}

A formula over weight pairs is \emph{smooth set-multilinear} if it has
scopes with the properties listed in Lemma~\ref{lem:port-smoothness}:
every nonzero monomial at a gate contains exactly one weight from each
pair in that gate's scope and none from the other pairs, the children
of a sum have the same scope, and the children of a product have
disjoint scopes whose union is the parent's scope. By
Lemma~\ref{lem:port-smoothness}, $W$ is such a formula with output
scope all $N$ pairs. The \emph{witness inequality} is the conclusion
$W(\mathbf1)\ge f(a)$ of Lemma~\ref{lem:ef-negative-descent} below.

\begin{lemma}[A polynomial equational proof of splitting one pair]
\label{lem:pair-shannon}
Let $P$ be a smooth set-multilinear formula and let pair $i$ belong
to its output scope. After arbitrary substitutions at the other
pairs,
\[
 P|_{(w_{i,0},w_{i,1})=(1,1)}
 =P|_{(w_{i,0},w_{i,1})=(1,0)}
  +P|_{(w_{i,0},w_{i,1})=(0,1)}.
 \tag{\thesection.12}\label{eq:pair-shannon}
\]
There is an arithmetic equational derivation of polynomial size in
the formula size.
\end{lemma}
\begin{proof}
Induct over gates whose scope contains $i$. At a weight leaf, the
identity is $1=1+0$ or $1=0+1$; a zero leaf gives $0=0+0$.
At a sum, both children contain $i$; add their inductive identities
and reassociate. At a product, exactly one child contains $i$.
Apply its identity and distribute over the other child. The other
child is unchanged and is not doubled. Each gate contributes a
bounded number of ring identities, and their substituted formulas
have size at most a polynomial in the original formula size.
\end{proof}

\begin{lemma}[Negative descent gives an EF witness inequality]
\label{lem:ef-negative-descent}
Let $W$ be a size-$S$ smooth set-multilinear integer formula on $N$
weight pairs. Let $f$ be Boolean-valued. Suppose there are EF proofs
of $W(\operatorname{pin}(b))=f(b)$ for an arbitrary Boolean vector
$b$. Then there are EF proofs of
\[
 W(\mathbf1)\ge f(a)
\]
for an arbitrary Boolean vector $a$, with polynomial overhead
in $S$, $N$, the input proof lengths, the
size of a Boolean circuit for $f$, and the arithmetic bit length.
If $W(\mathbf1)=0$ also has an EF proof, this gives refutations of $f(a)$
with polynomial overhead in the same parameters.
\end{lemma}
\begin{proof}
Assume $W(\mathbf1)<f(a)$ and define
\[
 Q_a(\mathbf w)=W(\mathbf w)-f(a)M_a(\mathbf w).
\]
It has the same full scope as $W$ and $Q_a(\mathbf1)<0$.
The monomial with a variable Boolean subscript means
$M_a=\prod_i((1-a_i)w_{i,0}+a_iw_{i,1})$; this is smooth
set-multilinear in the weight pairs, with $a$ treated as parameters.
We describe an EF construction of a Boolean vector $b$ with
$Q_a(\operatorname{pin}(b))<0$.

Initially the prefix is empty. If $b_1,\ldots,b_{j-1}$ have been
chosen, let $q_{j-1}$ be the value of $Q_a$ with those pairs pinned
and all remaining pairs set to $(1,1)$. The induction assertion is
$q_{j-1}<0$. Let $u_j$ and $v_j$ be the two values obtained by
pinning the next pair to zero and to one. Lemma~\ref{lem:pair-shannon}
gives $q_{j-1}=u_j+v_j$. Define the extension bit $b_j$ to be zero
when $u_j<0$, and one otherwise. In the first case the selected
value is negative. In the second case $u_j\ge0$ and
$u_j+v_j<0$, so $v_j<0$. This proves $q_j<0$.

After $N$ steps, the pinning identity gives
\[
 0>q_N=Q_a(\operatorname{pin}(b))=f(b)-f(a)[a=b].
\]
If $a=b$, equal inputs to the circuit for $f$ give equal outputs,
so $f(b)=f(a)$ and the last expression is zero. If $a\ne b$,
the expression is $f(b)$, which is nonnegative. This is a
contradiction in both cases.

For the size assertion, introduce extension variables for the $N$
selector bits, for the arithmetic gates in the $2N$ substituted
evaluations, and for their comparisons. Every evaluation is a
substitution instance of a formula of size polynomial in $S+N$.
Use one instance of~\eqref{eq:pair-shannon} per step and the
binary integer implication
$u+v<0\wedge u\ge0\Rightarrow v<0$.
These are polynomial-size arithmetic proof instances. Equal inputs
give equal outputs of the circuit for $f$ by induction over its
gates. There are $N$ steps and hence only polynomially many such
instances. Extension variables share their computations rather than
expanding all prefixes. No table of $2^N$ assignments is written.
\end{proof}

\paragraph{Why this is an EF argument.}
The extension definition of $b_j$ depends on earlier selector bits.
If an arithmetic evaluation has circuit depth $d$, the direct selector
construction has dependency depth $O(Nd)$. Unfolding those definitions
can cost an exponential in that depth. The proof-size bound in
Lemma~\ref{lem:ef-negative-descent} supplies no subexponential
ordinary-Frege bound. Establishing such a bound requires an additional
argument.

\subsection{The refutation}

\begin{theorem}[Relative to the supplied contraction theorem]
\label{thm:direct-subexp-ef}
Assume the integral contraction compiler and its index bounds from
the supplied manuscript. Then every
unsatisfiable 3-CNF with $n$ variables and $m$ clauses has an Extended
Frege refutation of size
\[
 2^{O((n+m)/\sqrt{\log(n+m+2)})}.
\]
In particular, every fixed-density family has refutations of size
$2^{O(n/\sqrt{\log n})}$.
\end{theorem}
\begin{proof}
Construct the port-weighted network and its compiled formula $W$.
Lemma~\ref{lem:port-smoothness} gives its scope structure and size.
Corollary~\ref{cor:port-pinning-identity}, with the quantitative
derivation bound in Lemma~\ref{lem:rankone-derivation-size} and the
arithmetic-to-EF translation of Lemma~\ref{lem:integer-derivation-ef},
gives an EF proof of its pinning identity within the stated size. The compiler's all-one evaluation is zero because
the input is unsatisfiable. Its numerical bit computation, with all
values fixed, has a polynomial-size EF verification in the compiled
formula size. Thus the two premises of
Lemma~\ref{lem:ef-negative-descent} have proofs of the required size.

The original assignment $x$ induces a port assignment $a(x)$ by
copying the associated original variable to every occurrence in its
equality tree and to the ports introduced by subdivision. Every
equality predicate and pair-consistency predicate is then true, and
the clause predicates are the clauses of the input. Hence the input
3-CNF implies $f_{\rm net}(a(x))=1$, with a polynomial-size proof
obtained from these local identities. Apply
Lemma~\ref{lem:ef-negative-descent} after substituting $a(x)$.
The polynomial overhead preserves the displayed exponential bound.
\end{proof}

\subsection{Small examples}

\paragraph{Hand checks.}
For one edge and two factorized vertices, the base identity is
$\sum_\lambda c_ub_{uv}(\lambda)c_vb_{vu}(\lambda)
=c_uc_vB_{uv}$, with each scalar pulled out once. For a one-bit
crossing at views $(0,1,0,1)$, the diagonal test is zero, while the
one-bit crossing identity has the two signed contributions $1$ and
$-1$; the factorized proof preserves their cancellation. At views
$(0,1,1,0)$ the two contributions are $0$ and $1$ and the diagonal
test is one. These are two pinned $2\times2$-grid checks.

For the negative search alone, take
\[
 Q=w_{1,0}w_{2,0}-2w_{1,0}w_{2,1}
      +w_{1,1}w_{2,0}-w_{1,1}w_{2,1}.
\]
Its all-one value is $-1$. The first zero child has value $-1$,
so the first selected bit is zero. The next zero child has value
$1$, so the second selected bit is one; the final value is $-2$.
This checks both branches of the selector rule.

Finally consider the four clauses
$(x\lor y),(x\lor\neg y),(\neg x\lor y),(\neg x\lor\neg y)$.
At the four assignments $00,01,10,11$, respectively the first,
second, third, and fourth clause is false. Their consistent port
extensions have $f_{\rm net}=0$. An inconsistent port assignment
fails a vertex equality or an edge-consistency test. Thus the
pinning predicate is zero on every port assignment, and the
all-one count is zero, as required by this instance of the theorem.


\section{Size and bit width of the pinning derivation}
\label{sec:rankone-size-audit}

This section proves the quantitative bound for the edge-factorized
pinning derivation, Lemma~\ref{lem:rankone-derivation-size}, which
the proof of Theorem~\ref{thm:direct-subexp-ef} uses. It assumes the
stated contraction compiler and the explicit local merge formulas
from the supplied manuscript. The conclusion concerns an arithmetic
equational derivation with shared definitions. Converting this
derivation to Extended Frege uses the separate binary-arithmetic
translation, Lemma~\ref{lem:integer-derivation-ef}. No conversion to
ordinary Frege is asserted here.

Lemma~\ref{lem:rankone-local-domains} bounds the local label domains
and the number of histories,
Lemma~\ref{lem:rankone-factor-formulas} bounds the formulas for the
auxiliary factors, and Lemma~\ref{lem:rankone-derivation-size}
assembles the derivation and bounds its bit widths.

Write $s$ for the network parameter, let $h_*$ be the initial
name length, put $q_*=2^{h_*}$, and set
\[
 L=\frac{s}{\sqrt{h_*}}.
\]
The compiler chooses $h_*=\Theta(\log s)$ and $q_*^3\le s$.
At a clique of name length $h$ and $q=2^h$ nodes, its potential
satisfies
\[
 \Phi\le C L^2h,
 \qquad
 \Phi=(q-1)^2\sum_{u<v}a_{uv}.
 \tag{\(*\)}
\]
All constants below may depend on the fixed compiler. Its label
code for budget $a$ has length at most $30\sqrt a$, with the
empty code for $a=0$.

\begin{lemma}[All local tables fit in the target exponent]
\label{lem:rankone-local-domains}
For every clique occurring in the recursion, the argument of one
entire node table has code length $O(L)$. The arguments comprising
the four old views and one message at one crossing also have
combined code length $O(L)$. The total number of clique-history
prefixes is $2^{o(L)}$.
\end{lemma}
\begin{proof}
At node $v$, Cauchy--Schwarz and $(*)$ give
\begin{align*}
 \sum_{u\ne v}30\sqrt{a_{uv}}
 &\le30\sqrt{(q-1)\sum_{u\ne v}a_{uv}}\\
 &\le30\sqrt{\frac{\Phi}{q-1}}
 \le30\sqrt C\,L\sqrt{\frac{h}{2^h-1}}=O(L).
\end{align*}
The last ratio is bounded for all integers $h\ge1$, since
$2^h-1\ge h$. This bound includes all incident labels, not just
one selected edge label. Invalid binary strings can also be
enumerated within the same bound and assigned value zero.

For a rectangle, let $a,b$ be its diagonal budgets and put
$H=a+b$. There are two old copies for each diagonal. The total
old-view code length is at most $60(\sqrt a+\sqrt b)$.
The message budget in either orientation is at most $2H$,
because $|k|\le2H$. Its code length is at most
$30\sqrt{2H}$. Since
$H\le\Phi/(q-1)^2$, these lengths together are at most
\[
 O\!\left(\frac{\sqrt\Phi}{q-1}\right)
 =O\!\left(L\frac{\sqrt h}{2^h-1}\right)=O(L).
\]

At a merge there are fewer than $q^2$ rectangles, and one
grade and one split field of $O(\log s)$ bits per rectangle.
Over all name lengths there are at most $O(\log h_*)$ merges.
Thus the total number of bits in a complete history is at most
\[
 D=O(q_*^2\log s\log h_*)
   =O(s^{2/3}\log s\log\log s)=o(L).
\]
There are at most $(1+\log_2h_*)2^D$ prefixes, including
prefixes of different lengths. Its base-two logarithm is $o(L)$.
\end{proof}

\begin{lemma}[Auxiliary factors have small formula representations]
\label{lem:rankone-factor-formulas}
In the edge-factorized recurrence
\[
 T_v=c_v\prod_{u\ne v}b_{vu}(\theta_{vu}),
\]
every particular value of every scalar $c_v$, every endpoint
factor $b_{vu}$, every edge inner product
$B_{uv}=\sum_\theta b_{uv}(\theta)b_{vu}(\theta)$,
and every local weighted crossing contraction has an arithmetic
formula of size $2^{O(L)}$. This holds uniformly for every
fixed history and every fixed label argument. The nonconstant
leaves are the original Boolean pin bits or bounded-arity
Boolean predicates of them.
\end{lemma}
\begin{proof}
The initial factors are Boolean predicates of the pin bits and
equalities between a fixed encoded label and the endpoint's pinned
bit string. They have size polynomial in $s$.

Consider a merge whose old name length is $h$. A new endpoint
factor is a product of at most $g+\binom g2=O(q)$ factors,
where $g=2^{h/2}$ and $q=g^2$. A vertical-stream factor is an
old endpoint factor at a fixed label. A rectangle factor is a sum
over the two old views held at that endpoint, of their two old
endpoint factors multiplied by a fixed routing coefficient.
The summation index therefore has length at most
\[
 O(\sqrt H)
 =O\!\left(L\frac{\sqrt h}{2^h-1}\right).
 \tag{\thesection.1}\label{eq:rankone-index-length}
\]
Invalid-packing and stream-domain tests are fixed scalar factors.
On any path through the recursively expanded factor formula,
there is at most one such new summation for each old name length.
The product above adds $O(\log q)=O(h)$ to the logarithm of
the fan-in on that path.

For the scalar recurrence, $c'_A$ is a product of the $g$ old
scalars and the inner products for all horizontal pairs inside the
row. There are fewer than $g^3$ such scalar or pair factors, so
the logarithm of the product fan-in is $O(h)$. A path either
continues into one old scalar, or enters one edge inner product
and thereafter follows endpoint-factor recurrences. The latter
inner product introduces only one label sum, bounded
by~\eqref{eq:rankone-index-length}.
Thus a scalar path does not acquire a separate sequence of
scalar recurrences after it has entered an endpoint factor.

The name lengths are powers of two. The sum of their values is
at most $2h_*=O(\log s)$. The sum of
$\sqrt h/(2^h-1)$ over this sequence is bounded by its sum
over all positive integers, which converges: for $h\ge2$,
the summand is at most $2\sqrt h/2^h$, and the latter
series is dominated eventually by a geometric series.
Consequently the sum of all index lengths~\eqref{eq:rankone-index-length}
on a path is
$O(L)$, while the product-fan-in contributions sum to
$O(\log s)=O(L)$.

A path may instead end inside a numerical routing-coefficient
formula. The compiler's transport bound contributes
$O(\sqrt\tau+\log^2s)$ bits, where $\tau$ is the parameter of that
bound (written $W$ in the supplied manuscript); the message tests
ensure $\tau\le3H$. This adds $O(L)$ by~\eqref{eq:rankone-index-length} and
$\log^2s=o(L)$. It does not lead back into an endpoint-factor
or scalar recurrence. A local weighted crossing uses four old
views, so it adds only a constant multiple
of~\eqref{eq:rankone-index-length} before following
one of the already bounded paths.

These facts bound by $O(L)$ the sum of the logarithms of the
fan-ins on any root-to-leaf path of the expanded formula; the
number of gates on a path is polynomial in $s$. To convert this
to a size bound, choose at every gate a child uniformly. The
probability of a particular leaf is the reciprocal of the product
of the fan-ins on its path, and hence is at least $2^{-O(L)}$.
These probabilities sum to one, so there are at most $2^{O(L)}$
leaves. The total gate count is at most the leaf count times the
maximum path length, which remains $2^{O(L)}$.
\end{proof}

\begin{lemma}[Size and bit width of the pinning derivation]
\label{lem:rankone-derivation-size}
For all-port pinning after the initial-clique construction, the
edge-factorized induction proving the compiler's output equal to
the fully pinned Boolean network predicate has an arithmetic
equational derivation with shared definitions of total size
$2^{O(L)}$. Every expression and every numerical intermediate
in this derivation can be represented with $2^{O(L)}$ signed
binary bits, uniformly for all Boolean pin assignments. The
derivation uses ring identities, bounded-arity Boolean identities
and Boolean circuit-defining equations for the initial pins, the
explicit packing bijections, and scalar local identities checked
by finite arithmetic computations. The Boolean identities have
constant-size truth-table proofs before substitution, so their
Extended Frege proofs add only polynomial overhead in the written
expressions. It never enumerates all labels of an entire large
clique simultaneously.
\end{lemma}
\begin{proof}
First create definitions for every endpoint-factor value, scalar,
edge inner product, and original table-entry value needed for
every history prefix. There are $2^{o(L)}$ prefixes. Each has
polynomially many nodes and edges. Each full node-table argument
has $O(L)$ bits by Lemma~\ref{lem:rankone-local-domains}.
Every value therefore belongs to an explicitly enumerable set of
size $2^{O(L)}$. Its formula has size $2^{O(L)}$:
for the auxiliary factors this is
Lemma~\ref{lem:rankone-factor-formulas}; for original table
entries it is the compiler's root-to-leaf index bound with the
global history fixed. Multiplying these bounds still gives
$2^{O(L)}$ total definitions.

Prove the table-factorization identities forward through the
history levels. For one table argument, expand its merged-table
sum over the old copies in one row or column. The compiler bounds
this index length by $O(L)$. Substitute the already established
old table-factorization identities. The old copies partition into
the disjoint horizontal-pair, vertical-stream and rectangle sets.
Regrouping factors and distributing these finite sums yields the
new scalar and endpoint factors. All expansion terms number
$2^{O(L)}$; each contains polynomially many factor occurrences.
An explicit sequence of distributive-law applications can factor
one independent index at a time. At every stage the number of
displayed terms is at most the original number of full-index
settings times a polynomial in the number of factors. Reordering
factors by adjacent swaps requires at most the square of that
number of factors per term. Hence the complete equational
derivation for this argument has size $2^{O(L)}$.
Doing this for all arguments and histories preserves the bound.

Next prove clique-value identities backward from the two-node
base cases. At a base case, the value formula is one label sum of
the product of the two tables. The table identities immediately
give the two scalars times their edge inner product. At an
earlier clique, use the already established identity for every
child clique once, referring to its shared definition and proof.
Its edge inner products can be unpacked and factored separately,
one new edge at a time. Each such label domain has $O(L)$ bits.
This step does not expand the product into simultaneous choices
of labels on all new edges.

For a fixed grade list, the orientation is fixed before the split
sum. Distributing only over splits gives for each rectangle the
scalar contraction of its two half-pieces. The relevant local
view and message domain has $O(L)$ bits. The row and column
factorizations both reconstruct the same weighted piece after
the entire split sum, so the orientation disappears at this
point. Finally, the complete grade lists have at most $o(L)$
bits in total. Distributing that grade sum across the rectangle
factors is therefore permissible within $2^{O(L)}$ size.
Apply the local piece-sum identity separately at each rectangle.
This produces the two old diagonal inner products, completing
the clique identity.

All numerical local identities can, if desired, be certified by
complete tables. For every fixed local view tuple, enumerate its
allowed grades, splits and messages, evaluate the transport
formulas, and check the asserted integer equality. There are
$2^{O(L)}$ such local tuples over all histories and each
evaluation has formula size $2^{O(L)}$. The packing tests and
bijections are likewise over one local domain of $O(L)$ bits.
Thus the arithmetic transcripts add $2^{O(L)}$ size. The
validity of the scalar equalities is guaranteed by the supplied
local identities, while the transcripts supply explicit finite
checks rather than leaving their universal formalization as an
additional hypothesis.

It remains to bound bit widths; arithmetic circuit size alone
would not do so in the presence of repeated squaring. Every
individual expression just used has an expanded formula of size
$T=2^{O(L)}$. For auxiliary quantities this was proved above
by a path bound, not inferred from circuit size. Local expansions
and substitutions multiply this size by at most another
$2^{O(L)}$ factor. Primitive leaves have polynomially many
bits (in fact the input predicates are Boolean and the transport
leaves are polynomially bounded integers). If each leaf has at
most $B_0$ bits, the inequalities
\[
 \log_2(|x+y|+1)\le1+\max\{\log_2(|x|+1),\log_2(|y|+1)\},
 \qquad
 \log_2(|xy|+1)\le
      \log_2(|x|+1)+\log_2(|y|+1)
\]
show by induction on the formula that all values have at most
$O(T(B_0+1))=2^{O(L)}$ bits. This is uniform over the Boolean
pin bits. More strongly, Lemma~\ref{lem:formula-integer-width}
applied to these expanded formulas bounds their adaptive full-result
implementation widths by $2^{O(L)}$. Choose one upper bound $B$
for all these widths, including the local rewrite instances and
their auxiliary expressions. Each operation uses its own full-result
width, at most $B$; it does not multiply two already $B$-bit words
and truncate their product back to $B$ bits. Equality checks
sign-extend to a common width. Thus all intermediate computations
represent exact integers without modular wraparound.
\end{proof}

\paragraph{A check on one $2\times2$ merge.}
Let the old nodes be $x,y$ in the upper row and $u,v$ in the
lower row. Assume each table has the form
$c_t\prod_{t'\ne t}b_{tt'}$. In the row orientation, the
upper scalar is $c_xc_yB_{xy}$ and the lower scalar is
$c_uc_vB_{uv}$. The new pair label has three streams: the
vertical labels for $xu$ and $yv$, and one rectangle message.
Its upper endpoint factor is
$b_{xu}(\lambda_1)b_{yv}(\lambda_2)G_{\rm up}(m)$,
and its lower factor is
$b_{ux}(\lambda_1)b_{vy}(\lambda_2)G_{\rm down}(m)$.
The inner product separates into
$B_{xu}B_{yv}\sum_mG_{\rm up}(m)G_{\rm down}(m)$.
Summing the rectangle split and grade changes the last factor
to $B_{xv}B_{yu}$. Thus the result is
\[
 c_xc_yc_uc_v
 B_{xy}B_{uv}B_{xu}B_{yv}B_{xv}B_{yu},
\]
one scalar per node and one inner product for each of the six
old pairs. This calculation also identifies why the whole split
sum must precede the local piece identity.



\section{Removing the clause-density restriction}
\label{sec:ef-sparsification}

This section transfers the bound of
Theorem~\ref{thm:direct-subexp-ef} to all normalized 3-CNFs, using the
IPZ sparsification tree, and proves
Theorem~\ref{thm:all-3cnf-ef}. Its purpose is to justify the
variable-only statement without silently choosing a density-dependent
rate.

The source is Impagliazzo, Paturi, and Zane~\cite{ipz}, \S2:
Theorem~1, printed p.~520; \textsc{Reduce}, p.~521;
Lemma~7 and Corollary~1, p.~523. For every fixed $\varepsilon>0$,
its binary tree represents a width-three CNF $F$ as a disjunction
of at most $2^{\varepsilon n}$ leaves, each with at most
$C_\varepsilon n$ clauses. One may use the universe of $2n$
literals and replace the set-system parameter by $\varepsilon/2$.

\begin{lemma}[A sparsification split has a polynomial Frege proof]
\label{lem:ef-sparse-split}
At every node of this tree, if $A$ is the parent conjunction and
$B,C$ are its child conjunctions, the implication
$A\to(B\lor C)$ has an ordinary Frege proof of length
polynomial in $n$, with a polynomial independent of the
sparsification thresholds.
\end{lemma}
\begin{proof}
A split selects clauses $H\lor P_1,\ldots,H\lor P_r$.
One child adds the heart $H$; the other adds all petals $P_i$.
Both children may delete clauses that are supersets of retained
clauses. Every clause of $B$ is either inherited from $A$ or is
$H$. Thus projection from conjunctions and conjunction introduction
give $(A\land H)\to B$.

Under $A\land\neg H$, each selected clause forces its petal
$P_i$. Projection of the selected clause followed by the fixed
tautology
\[
 (A\to(H\lor P_i))\to((A\land\neg H)\to P_i)
\]
proves this implication for every petal. Inherited clauses also
follow from $A$. Combine them to obtain
$(A\land\neg H)\to C$. The fixed two-case tautology then
gives $A\to(B\lor C)$.

Deleting supersets causes no difficulty: a retained clause
contained in a deleted clause implies it by weakening. All
clauses have width at most three. Every node uses only the
$O(n^3)$ distinct clauses on the same literal universe, so its
written length is polynomial in $n$, including variable names.
The construction above uses at most polynomially many
projections, weakenings, and conjunction-introduction steps.
Each is a substitution instance of a fixed tautology. Substitution
through a fixed proof multiplies its size only by the written
substitution length. This proves the uniform polynomial bound.
\end{proof}

\begin{lemma}[Leaf refutations recombine additively]
\label{lem:ef-sparse-combine}
If the leaves have EF proofs of their negated conjunctions of
sizes $P_1,\ldots,P_r$, the root has an EF proof of its negated
conjunction of size
\[
 \sum_{i=1}^{r}P_i+r\,n^{O(1)}.
\]
\end{lemma}
\begin{proof}
Rename the extension variables of the leaf proofs so that they are
mutually disjoint and disjoint from the original variables. Write
each leaf proof once. Process the tree in postorder. At a parent,
the two child negations have already been derived. Combine them
with its split implication using the fixed tautology
\[
 (A\to(B\lor C))\to(\neg B\to(\neg C\to\neg A)).
\]
No previously written proof is copied at this step; its concluding
line is reused. A full binary tree with $r$ leaves has $r-1$
internal nodes: if it has $I$ internal nodes, counting child edges
gives $2I$, and every one of the $I+r-1$ nonroot vertices is the
lower end of exactly one child edge, so $2I=I+r-1$. The preceding
lemma therefore gives the displayed additive bound.
\end{proof}

\paragraph{Converting clause refutations before recombination.}
If a leaf conjunction $A$ instead has a refutation from its clauses
of written size $R$, first convert it into a proof of $\neg A$.
Rename all extension variables to be fresh and disjoint from the
variables of $A$. Keep every extension definition as an unconditional
extension axiom, in its original topological order. For an original
clause axiom $D$ of the leaf, derive $A\to D$ by conjunction
projection. For an extension axiom or an ordinary Frege axiom $D$,
use that axiom and the fixed weakening tautology to derive $A\to D$.
Conditionalize each Frege inference using its fixed propositional
template; for modus ponens this is
\[
 (A\to D)\to((A\to(D\to E))\to(A\to E)).
\]
The final line is $A\to\bot$, equivalent to $\neg A$.
The extension definitions remain legal because no fresh variable is
introduced into an earlier definition or into $A$. A bounded amount
of proof work is added per original line, with occurrences of $A$
of length $n^{O(1)}$, so the converted size is
$n^{O(1)}(R+1)$. Apply the additive lemma to these converted proof
sizes, not to the unconverted clause-refutation sizes.

\begin{theorem}[Variable-only Extended Frege upper bound]
\label{thm:all-3cnf-ef}
Relative to the supplied integral contraction construction, every
unsatisfiable normalized $n$-variable 3-CNF has an Extended Frege
refutation of size $2^{o(n)}$.
\end{theorem}
\begin{proof}
Let $p(n)$ be the maximum, over such unsatisfiable formulas, of
their shortest EF refutation sizes, and let $p^{\neg}(n)$ be the
maximum, over the same formulas $F$, of the size of a shortest EF
proof of the negated conjunction $\neg F$; both are at least one. Fix
$\varepsilon>0$ and use its IPZ tree. Every leaf has at most
$C_\varepsilon n$ clauses and is unsatisfiable: any satisfying
leaf assignment would satisfy the parent, since the child only
adds stronger clauses before deleting redundant supersets.
Theorem~\ref{thm:direct-subexp-ef}, followed by the deduction
conversion above, gives proofs of the leaf negations of size
\[
 2^{K_\varepsilon n/\sqrt{\log(n+2)}}
\]
for a constant $K_\varepsilon$. Lemma~\ref{lem:ef-sparse-combine}
implies
\[
 p^{\neg}(n)\le
 2^{\varepsilon n}
 \left(2^{K_\varepsilon n/\sqrt{\log(n+2)}}+n^{O(1)}\right).
\]
A proof of $\neg F$ gives a refutation from the clauses of $F$:
conjunction introduction derives $F$ from its clauses in size
$n^{O(1)}$, and the fixed tautology $F\to(\neg F\to\bot)$ with two
applications of modus ponens derives $\bot$. Hence
\[
 p(n)\le p^{\neg}(n)+n^{O(1)},
\]
where the polynomial does not depend on $\varepsilon$. Combine the
two displayed bounds, divide the base-two logarithm by $n$, and let
$n$ tend to infinity. For this fixed $\varepsilon$, the limit superior is
at most $\varepsilon$. This holds for every positive
$\varepsilon$, while the ratio is nonnegative. Its limit is
therefore zero, which is exactly $p(n)=2^{o(n)}$.
No uniform estimate in $C_\varepsilon$ and no unspecified
choice of $\varepsilon(n)$ is used.
\end{proof}

\begin{corollary}[Every fixed clause width]\label{cor:fixed-width}
Relative to the same contraction construction, for each fixed $k$ the
maximum shortest EF refutation size of an unsatisfiable normalized
$n$-variable CNF of width at most $k$ is $2^{o(n)}$.
\end{corollary}
\begin{proof}
Fix $k$ and $\varepsilon>0$. Apply the width-$k$ IPZ tree
\cite[Theorem~1 and Corollary~1]{ipz} before introducing any variables.
It has at most $2^{\varepsilon n}$ leaves, each with at most
$C_{k,\varepsilon}n$ clauses. The split proof of
Lemma~\ref{lem:ef-sparse-split} works verbatim for the same heart
and petal implications. Each node now contains at most $O_k(n^k)$
distinct clauses, so each split costs $n^{O_k(1)}$ symbols.

In a leaf, encode each clause of length greater than three by a
binary OR tree. Give every internal gate a fresh extension variable
and write the three defining clauses for $y\leftrightarrow(a\lor b)$.
Assert the root. From the original clause and the extension axioms,
EF derives the defining clauses and the root unit in size polynomial
in the leaf length: at fixed $k$ this is a fixed propositional
derivation up to renaming and signs of literals. Short clauses are
retained. The resulting unsatisfiable 3-CNF has
$O_{k,\varepsilon}(n)$ variables and clauses, so
Theorem~\ref{thm:direct-subexp-ef} refutes it in size
$2^{O_{k,\varepsilon}(n/\sqrt{\log(n+2)})}$. Append that refutation
with fresh extension names, using the derived clauses as its premises.
This gives a refutation from the original leaf clauses.

Convert each leaf refutation to a proof of its negation as above and
apply Lemma~\ref{lem:ef-sparse-combine}, with $n^{O_k(1)}$ per node.
The base-two logarithm of the root proof size divided by $n$ has
limit superior at most $\varepsilon$. Taking the infimum over fixed
$\varepsilon>0$ proves the claim, including at density $\Theta(n^{k-1})$.
\end{proof}

\begin{theorem}[Quantitative sparsification in EF]\label{thm:quant-width}
Under the same contraction premise, every normalized unsatisfiable
$n$-variable CNF of fixed width $k\ge3$ has an EF refutation of size
\[
 2^{\,O_k\left(n(\log\log n)^{1-1/k}/(\log n)^{1/(2k)}\right)} .
\]
At width three the exponent is
$O(n(\log\log n)^{2/3}/(\log n)^{1/6})$.
For growing $k\ge3$ satisfying $k\log(k+2)=o(\log\log n)$, the size
is $2^{o(n)}$. The proofs can be constructed within the respective
bounds, with the input-reading cost added for redundant encodings.
\end{theorem}
\begin{proof}
Use Calabro's constructive sparsification algorithm~\cite{calabro}.
For $0<\epsilon\le1$ it gives at most $2^{\epsilon n}$ leaves with
$n^{O(k)}$ delay. The penultimate bound in equation~(5.4) of that
source, with the subclause size set to one, bounds the occurrences
of each literal by
\[
 \Delta=\big[C k^2\epsilon^{-1}\log(k/\epsilon)\big]^{k-1}
\]
for an absolute $C$. The children add a heart or its petals and then
discard supersets, exactly the implication in Lemma~\ref{lem:ef-sparse-split}.
Every intermediate conjunction and split proof has size $n^{O(k)}$.
Each leaf has $O(n\Delta)$ clauses. The OR-tree construction of
Corollary~\ref{cor:fixed-width} uses $O(kn\Delta)$ variables and
clauses. Uniformly in $k$, its root unit follows by propagating each
original literal up its path to the root and eliminating the literals
from the original clause. This needs only polynomial work in the
written tree and clause.

Use Corollary~\ref{cor:ef-proof-output} to construct each refutation.
The leaf time exponent is $O(kn\Delta/\sqrt{\log n})$, by monotonicity
of $v/\sqrt{\log(v+2)}$ at large $v$. The deduction conversion and
bottom-up reuse of proofs from Lemma~\ref{lem:ef-sparse-combine} add
at most polynomial overhead in each leaf proof and $n^{O(k)}$ per
tree node. Variable names can include a leaf index of $O(\epsilon n+1)$
bits. The base-two logarithm of both total size and construction time
is at most
\[
 O\big(\epsilon n+kn\Delta/\sqrt{\log n}+k\log n\big).
\]
For fixed $k$ take $\epsilon=(\log\log n)^{1-1/k}/(\log n)^{1/(2k)}$.
Then $\log(k/\epsilon)=O_k(\log\log n)$ and
$\Delta/\sqrt{\log n}=O_k(\epsilon)$. For growing $k$, put
$L=\log\log n$ and take $\epsilon=2^{-\sqrt{L/k}}$. The width
condition gives $\epsilon=o(1)$ and
$\log\Delta=O(k\log(k+2)+\sqrt{kL})=o(L)$. Hence
$k\Delta=(\log n)^{o(1)}$ and all three terms are $o(n)$.
Dyadic rounding within a factor two preserves each estimate. The
width is read directly from a normalized formula.
\end{proof}

\paragraph{Comparison on hard families.}
At sufficiently large fixed clause density, random 3-CNFs are
unsatisfiable with high probability. On the same distribution, exponential
lower bounds hold for Resolution length~\cite{BSW}, for polynomial
calculus monomial size, for example over $\mathbb F_3$~\cite{BI},
and for SOS monomial size with the clause-satisfaction equations
$p_C(x)=1$~\cite[Corollary~14]{AH}. Here $p_C$ is the multilinear
indicator of satisfying $C$. The fixed-density EF bound is
$2^{O(n/\sqrt{\log n})}$ on every unsatisfiable instance in that
distribution. This gives both a high-probability comparison and
worst-case families witnessing the gap. At other widths and densities
the EF upper bound still applies, while a size comparison requires
a lower bound in the corresponding system and encoding.

\paragraph{A four-clause split check.}
For
$A=(x\lor y)\land(x\lor\neg y)\land(\neg x\lor y)
\land(\neg x\lor\neg y)$, choose heart $x$ in the first
two clauses. One child adds $x$; the other adds $y$ and
$\neg y$. If $x=1$, the first child follows from $A$; if
$x=0$, the first two clauses force $y$ and $\neg y$, yielding
the second child. The latter is immediately contradictory. In
the first child the last two clauses force the same contradiction.
This checks both the split and the additive recombination.



\section{Exact counting certificates in CLIP plus Extended Frege}
\label{sec:clip-counting}

This section applies the pinning identity to an existing counting
proof system; it introduces no new proof format. The CLIP framework of
Chede, Chew, and Shukla~\cite{ccs} asks for a circuit computing
cumulative model counts and a propositional proof of its base,
successor, and final-value conditions. Their Definition~19 and
Lemma~21 show how short EF proofs of full restriction and one-variable
splitting give the required cumulative-count certificate; their
Theorem~23 states the corresponding p-simulation theorem. These
locators refer to revision~1 of ECCC TR24-081, PDF pp.~14--15.

Lemma~\ref{lem:ef-pin-equality} lifts equality at all pins to
equality at all-one weights, Lemma~\ref{lem:original-weight-pinning}
moves from port weights back to one weight pair per original variable,
and Theorem~\ref{thm:clip-counting} builds the certificate.

\begin{lemma}[Equality at all pins implies equality at all-one weights]
\label{lem:ef-pin-equality}
Let $P$ and $Q$ be smooth set-multilinear integer formulas on the
same $r$ weight pairs. Suppose an EF proof establishes
$P(\operatorname{pin}(b))=Q(\operatorname{pin}(b))$ for an arbitrary
Boolean vector $b$. Then EF proves $P(\mathbf1)=Q(\mathbf1)$ with
polynomial overhead in the formulas, the input proof, and the
arithmetic bit length.
\end{lemma}
\begin{proof}
Put $D=P-Q$. Its pinning identity is $D(\operatorname{pin}(b))=0$.
Apply Lemma~\ref{lem:ef-negative-descent} with the Boolean function
identically zero, first to $D$ and then to $-D$. It gives
$D(\mathbf1)\ge0$ and $-D(\mathbf1)\ge0$. Integer order then gives
$D(\mathbf1)=0$. Both applications have the asserted polynomial
overhead.
\end{proof}

\begin{lemma}[Returning to the original variables]
\label{lem:original-weight-pinning}
For a 3-CNF $F$ on its original variables $x_1,\ldots,x_n$, there is
a smooth set-multilinear formula $U$ with one weight pair per original
variable and size $2^{O(s/\sqrt{\log s})}$, where $s=O(n+m)$, for which
EF has a proof of
\[
 U(\operatorname{pin}(x))=F(x)
\]
of size $2^{O(s/\sqrt{\log s})}$. Also $U(\mathbf1)$ is the number
of satisfying assignments of $F$.
\end{lemma}
\begin{proof}
For every original variable, choose one representative port in its
connected equality network. Keep the weight pair at that port and
set all other port weights to one in the formula $W$ from
Section~\ref{sec:direct-ef}; call the resulting formula $U$.
Deleting fixed weight pairs from every declared scope preserves
smooth set-multilinearity. The formula size does not increase.
For a variable with no occurrence, the separate edge between two
degree-one predicates identically one was included before making
the initial clique. It supplies that variable's independent binary
choice and representative port at constant network cost.

Fix the representative ports to $x$ but leave the remaining port
weights free. Let this restricted formula be $P_x$. Each remaining
port $p$ belongs to one original variable; denote its index by
$i(p)$. Define
\[
 Q_x=F(x)\prod_{p\text{ not representative}}w_{p,x_{i(p)}}.
\]
The two formulas have the same scope. Under an arbitrary pinning
$b$ of the remaining ports, the pinning theorem for $W$ identifies
$P_x(\operatorname{pin}(b))$ with the network predicate whose
representative bits are $x$ and whose other port bits are $b$.

The equality predicates and pair-consistency tests connect exactly
the copies of each original variable. They hold if and only if every
such copy agrees with its representative. One direction follows by
substituting equal bits into each local predicate. For the other
direction, take a path in the equality tree from a port to its
representative and apply transitivity of equality once per edge.
All trees have total size $O(s)$, so proving these equivalences for
all ports takes polynomial size. Under these equalities, every clause
predicate becomes the corresponding clause of $F(x)$. Thus EF proves
\[
 P_x(\operatorname{pin}(b))
 =F(x)\prod_{p\text{ not representative}}[b_p=x_{i(p)}]
 =Q_x(\operatorname{pin}(b))
\]
with polynomial overhead in the all-port pinning proof.
Lemma~\ref{lem:ef-pin-equality} gives
$P_x(\mathbf1)=Q_x(\mathbf1)=F(x)$. The left side is
$U(\operatorname{pin}(x))$. At all-one weights the network count
is unchanged, so $U(\mathbf1)=\#F$.
\end{proof}

\begin{theorem}[Subexponential CLIP plus EF counting proofs]
\label{thm:clip-counting}
Relative to the supplied integral contraction compiler, every
3-CNF on $n$ variables with $m$ clauses has a CLIP plus Extended
Frege proof of its exact model count of size
\[
 2^{O((n+m)/\sqrt{\log(n+m+2)})}.
\]
\end{theorem}
\begin{proof}
Use the formula $U$ and pinning proof from
Lemma~\ref{lem:original-weight-pinning}. For a prefix
$a_1,\ldots,a_j$, write
\[
 C(a_1\cdots a_j)
 =U(\operatorname{pin}(a_1,\ldots,a_j),\mathbf1
       \text{ on the remaining pairs}).
\]
This is one substitution instance of $U$. The pinning theorem gives
$C(a_1\cdots a_n)=F(a)$, and
Lemma~\ref{lem:pair-shannon} gives
\[
 C(p)=C(p0)+C(p1)
 \tag{\thesection.1}\label{eq:clip-prefix-split}
\]
for every prefix length and arbitrary Boolean prefix $p$. These
are precisely the two properties of restriction required by
Chede--Chew--Shukla's Definition~19, and are the two identities used
in their Lemma~21. We instantiate the cumulative-count construction
explicitly below.

For completeness, one explicit cumulator is
\[
 H(a)=\sum_{i=1}^n a_i\,C(a_1\cdots a_{i-1}0)+F(a),
 \tag{\thesection.2}\label{eq:clip-cumulator}
\]
where $a_1$ is the most significant bit. Every summand counts the
assignments whose first difference from $a$ is a zero at a position
where $a$ has a one. These prefix cylinders are disjoint, and the
last term counts $a$ itself. The circuit uses $n$ copies of $U$
and polynomial-size binary arithmetic.

The inductive proposition can also be proved directly from
\eqref{eq:clip-prefix-split}. At $a=0^n$, all summands of
\eqref{eq:clip-cumulator} vanish and $H(a)=F(a)$.
For a successor pair, write
\[
 a=p0\,1^r,\qquad b=p1\,0^r.
\]
The contributions from positions in $p$ agree. Repeated application
of~\eqref{eq:clip-prefix-split} along the suffix of the zero child
shows that all the remaining contributions to $H(a)$ together equal
$C(p0)$. For $H(b)$, the new one at the displayed differing position
contributes $C(p0)$, every later bit is zero, and the final term is
$F(b)$. Thus $H(b)=H(a)+F(b)$. There are at most $n$ possible lengths
$r$; for each case the proof uses at most $n$ split identities.
The cases exhaust successor pairs by the bit definition of binary
increment. At $a=1^n$, the same telescoping identity gives
$H(a)=C(\varnothing)=U(\mathbf1)$, whose numerical evaluation is
the claimed model count.

To match CLIP's prescribed $n+1$ unsigned output bits, one may
initially perform the arithmetic with the larger exact bit width
used in Section~\ref{sec:direct-ef}. Each restricted count satisfies
\[
 0\le C(p)\le2^{n-|p|}
 \tag{\thesection.3}\label{eq:clip-count-range}
\]
with polynomial-size EF proofs. For the lower bound, apply
Lemma~\ref{lem:ef-negative-descent} to the restricted formula and
its Boolean pinning function. For the upper bound, apply it to
\[
 \prod_{i>|p|}(w_{i,0}+w_{i,1})-U|_{\operatorname{pin}(p)},
\]
whose pinning function is the Boolean complement of the same
restricted predicate. Its all-one value is
$2^{n-|p|}-C(p)$, so it is nonnegative. There are only $n$ prefix
lengths and the prefixes remain symbolic. Therefore these proofs
have polynomial total overhead.

Equations~\eqref{eq:clip-cumulator} and
\eqref{eq:clip-count-range}, together with $0\le F(a)\le1$, give
\[
 0\le H(a)\le\sum_{i=1}^n a_i2^{n-i}+1
            =\operatorname{num}(a)+1\le2^n.
\]
Hence the circuit's low $n+1$ bits equal the exact integer value,
with a short proof that no discarded high bit is nonzero. All
base, successor, and final-value equalities hold in the unsigned
encoding required by CLIP. The total circuit and proof length is
polynomial in the size budget $2^{O(s/\sqrt{\log s})}$.
\end{proof}

\paragraph{A two-variable check.}
For $F(x_1,x_2)=x_1\wedge\neg x_2$, the four lexicographic
assignments are $00,01,10,11$, and their cumulative counts are
$0,0,1,1$. Formula~\eqref{eq:clip-cumulator} gives those values:
$C(0)=0$, $C(10)=1$, and the final summand $F(a)$ contributes
one only at $10$. The transition $01\to10$ is the carry case
$p$ empty and $r=1$; it adds exactly $F(10)=1$.


\section{Consequences and the ordinary-Frege problem}
\label{sec:ef-boundaries}

\paragraph{What is proved.}
Theorems~\ref{thm:direct-subexp-ef} and
\ref{thm:all-3cnf-ef} give the claimed bounds in a fixed,
established propositional proof system. The construction does not
add a counting axiom to EF. Its long numerical identities are
verified computations, its symbolic transformations use elementary
integer laws with known EF proofs, and its nonnegativity argument
uses acyclic extension definitions and propositional reasoning.

\begin{corollary}[Proof output at the clause bound]
\label{cor:ef-proof-output}
For a 3-CNF with $n$ specified variables and $m$ clauses, the
construction can output an EF refutation, when the input is
unsatisfiable, within time
$2^{O((n+m)/\sqrt{\log(n+m+2)})}$. For active variables,
$n\le3m$, so the bound is $2^{O(m/\sqrt{\log(m+2)})}$.
\end{corollary}
\begin{proof}
The explicit compiler enumerates its formula, history records,
local tables and routing computations within the bound of
Lemma~\ref{lem:rankone-derivation-size}. That lemma also gives
explicit equational rewrites and numerical calculation traces of
the same asymptotic size. The arithmetic library constructs its
EF proof instances in time polynomial in their lengths. The
selector proof is obtained by writing a fixed template for each
port; its selector bits are variables in a proof, so the
constructor does not search over their values. The all-one
numerical evaluation determines whether the count is zero and,
when it is, supplies the ground computation used in the proof.
All these construction times are polynomial in
$2^{O((n+m)/\sqrt{\log(n+m+2)})}$, which preserves the
displayed bound.
\end{proof}

\paragraph{Extended Resolution.}
The same asymptotic upper bounds hold for Extended Resolution,
which polynomially simulates EF. One direct way to see the
simulation is to name every gate in every EF line by a fresh
extension variable and introduce its Tseitin definition. A fixed
Frege axiom schema on the finitely many named inputs has a
fixed resolution proof of its output under these definitions.
For modus ponens, the named circuits for $A$, $A\to B$, and
$B$ give a fixed local resolution derivation of the asserted
output of $B$ from the first two asserted outputs. Introducing
gate definitions and copying these fixed derivations for all
lines costs polynomially many clauses in the EF proof length.
For a tautology proof of $\neg F$, assume the clauses of $F$
and resolve its asserted negation to the empty clause. The
simulation is also recorded in Kraj\'i\v{c}ek~\cite{krajicek}, \S3.3,
printed p.~53, item~4 following Theorem~3.3.3. This gives a further
proof-system corollary; it does not improve the counting exponent.

\paragraph{Ordinary Frege is not settled.}
The original target statement concerned ordinary Frege and
total written proof size. EF may reuse named computations; Frege
must write formulas instead. The selectors in the negative
descent depend on earlier selectors. The arithmetic helper
circuits also require care if their names are eliminated.
A bound on EF proof size is not a bound on the sizes of the
resulting expanded formulas.

\begin{lemma}[A sufficient extension-depth bound]
\label{lem:extension-depth-transfer}
Let an EF proof have written size $P$, and suppose its extension
definitions have fan-in at most two and total dependency depth
at most $d$. If its conclusion contains no extension variables,
then it has an ordinary Frege proof of size $P\,2^{O(d)}$.
In particular, $P=2^{o(n)}$ and $d=o(n)$ suffice for a
$2^{o(n)}$-size ordinary Frege proof.
\end{lemma}
\begin{proof}
Expand each extension variable recursively into its defining
Boolean formula. A binary tree of depth $d$ has at most
$2^{d+1}-1$ nodes, proved by induction from the recurrence
$t(d)\le1+2t(d-1)$ with $t(0)=1$.
Consequently each variable occurrence in the proof expands to
at most $2^{O(d)}$ symbols. Apply this expansion to every
nonextension line. Frege axioms remain instances of their same
schemata, and modus ponens remains valid.

An extension line $z\leftrightarrow G$ becomes
$G^*\leftrightarrow G^*$, where $G^*$ is the fully expanded
definition. Replace that line by the fixed proof of
$A\leftrightarrow A$ with $A=G^*$. Its size is proportional
to $|G^*|$. The conclusion is unchanged because it has no
extension variables. Summing these costs over the original
written proof gives $P2^{O(d)}$. Taking base-two logarithms
proves the final assertion.
\end{proof}

\paragraph{The remaining quantitative target.}
The present proof does not establish the hypothesis $d=o(n)$.
With $N=O(n)$ ports for a fixed-density formula and an evaluator
of depth $d_0$, a direct selector implementation has a dependency
bound of order $N d_0$, before other helper definitions are
included. This is an upper-bound calculation, not a lower bound
on all implementations. A useful next attack is a proof of the
witness inequality in ordinary Frege that preserves the
subexponential size budget, either by controlling all necessary
formula expansions or by replacing the sequential selector
argument. The SOS obstruction of
Appendix~\ref{sec:sos-length-obstruction} explains why replacing the
signed polynomial by a short explicit squared norm cannot be
a universal solution.

\paragraph{What is not claimed.}
The results are an upper bound on refutation size and a
certificate-producing version of the counting simulation. It
does not prove $\mathrm P\ne\mathrm{NC}$, an unrestricted
Boolean circuit lower bound, or a separation between Frege
and Extended Frege.

\appendix



\section{An exact obstruction to explicit sum-of-squares positivity}
\label{sec:sos-length-obstruction}

\paragraph{Status and scope.}
This section proves a representation separation, conditional only on the
integral weighted formula compiler imported in Section~\ref{sec:imported}.
The lower-bound argument and the coding construction are unconditional.
The result does not prove a lower bound for Frege or Extended Frege.
Its examples are satisfiable instances: the obstruction concerns a
universal conversion of weighted-count formulas into explicitly listed
sums of polynomial squares.

Theorem~\ref{thm:sos-exact-criterion} characterizes when the
sum-of-squares length of an assignment polynomial is maximal.
Lemmas~\ref{lem:sos-minimal-code}--\ref{lem:sos-fast-minimal} supply
explicit codes meeting the criterion, and
Theorem~\ref{thm:sos-signed-separation} combines them with the
compiler. Lemma~\ref{lem:sos-weighted-descent} and
Corollary~\ref{cor:sos-ef-positivity} show that EF nevertheless proves
the corresponding positivity statements efficiently.

\paragraph{Definitions.}
For $a\in\{0,1\}^N$, introduce the assignment monomial
\[
 M_a(w)=\prod_{i=1}^{N}w_{i,a_i}.
\]
For nonempty $A\subseteq\{0,1\}^N$, put
\[
 P_A(w)=\sum_{a\in A}M_a(w)^2.
\]
The real polynomial sum-of-squares length, denoted
$\operatorname{sos}_{\mathbb R}(P_A)$, is the least $D$ such that
$P_A=\sum_{\ell=1}^{D}p_\ell^2$ for real polynomials $p_\ell$.
No restriction on the degree, coefficient signs, or computation size of
the $p_\ell$ is imposed in this definition.
We call $A$ \emph{pair-unique} when
\[
 M_aM_b=M_cM_d\quad\Longrightarrow\quad
 \{a,b\}=\{c,d\}
 \qquad(a,b,c,d\in A),
\]
where the pairs are unordered and may contain repeated elements.

\begin{lemma}[Forced block degree and support]
\label{lem:sos-forced-support}
If $P_A=\sum_{\ell=1}^{D}p_\ell^2$ over $\mathbb R$, then
\[
 p_\ell=\sum_{a\in A}v_{\ell,a}M_a
 \qquad(v_{\ell,a}\in\mathbb R).
\]
In particular, each square root is homogeneous of degree one in each
block $\{w_{i,0},w_{i,1}\}$.
\end{lemma}
\begin{proof}
Fix a block $i$, and let $e$ be the maximum degree in this block among
the nonzero $p_\ell$.  If $e>1$, the block-degree-$2e$ part of the
squared sum is the sum of the squares of the block-degree-$e$
homogeneous parts of the $p_\ell$.  This sum is zero, since $P_A$ has
block degree two.  A sum of squares of real polynomials is zero only
if each summand polynomial is zero: evaluation at any real point
gives a sum of nonnegative real numbers, and each must vanish.
A polynomial vanishing at every real point is the zero polynomial;
this follows by induction on the number of variables, viewing it as a
univariate polynomial and choosing a point where a nonzero coefficient
polynomial does not vanish.
We have contradicted the choice of $e$, so every $p_\ell$ has degree
at most one in the chosen block.

Set both variables of this block to zero.  The left side becomes zero,
so every specialized $p_\ell$ is zero by the same sum-of-squares
argument.  Its block-degree-zero part is therefore zero.
Every $p_\ell$ is homogeneous of degree one in this block.
Repeating this argument for all blocks gives an expansion
$p_\ell=\sum_{a\in\{0,1\}^N}v_{\ell,a}M_a$.

For $a\notin A$, set $w_{i,a_i}=1$ and $w_{i,1-a_i}=0$ for each $i$.
At this point $M_b$ is one if $b=a$ and zero otherwise.
Consequently $0=P_A=\sum_\ell v_{\ell,a}^2$, forcing all these
coefficients to be zero.
\end{proof}

\begin{theorem}[Exact maximal-length criterion]
\label{thm:sos-exact-criterion}
For every nonempty $A\subseteq\{0,1\}^N$,
\[
 \operatorname{sos}_{\mathbb R}(P_A)=|A|
 \quad\Longleftrightarrow\quad A\text{ is pair-unique}.
\]
When $A$ is pair-unique, the positive semidefinite Gram matrix of
$P_A$ in the assignment-monomial basis is uniquely the identity on
the coordinates in $A$, padded by zeros outside $A$.
When $A$ is not pair-unique,
$\operatorname{sos}_{\mathbb R}(P_A)\le |A|-2$.
\end{theorem}
\begin{proof}
Consider any sum-of-squares representation and use
Lemma~\ref{lem:sos-forced-support}.  For $a\in A$, write
$v_a=(v_{1,a},\ldots,v_{D,a})\in\mathbb R^D$.
The coefficient of $M_a^2$ gives $\langle v_a,v_a\rangle=1$.
Indeed, a product $M_bM_c$ can equal $M_a^2$ only when $b=c=a$:
in each block, the latter monomial contains only one of its two
variables, with exponent two.

Assume $A$ is pair-unique.  For distinct $a,b\in A$, the coefficient
of $M_aM_b$ is $2\langle v_a,v_b\rangle$, because no other unordered
pair contributes to that monomial.  This coefficient is zero in $P_A$.
Thus the $|A|$ vectors $v_a$ are orthonormal in $\mathbb R^D$, and
$D\ge |A|$.  The defining expression for $P_A$ uses $|A|$ squares.
The same coefficient calculation determines every entry of its Gram
matrix and proves the uniqueness assertion.

Conversely, suppose $M_aM_b=M_cM_d$ for two distinct unordered pairs.
The four words are distinct.  If a pair repeated one word, examination
of the exponents in every block would force the other pair to repeat
the same word.  If the pairs shared one word, cancellation of its
monomial would force them to share the other word as well.
Replace the four corresponding squares in the defining expression by
\[
 M_a^2+M_b^2+M_c^2+M_d^2
 =(M_a+M_b)^2+(M_c-M_d)^2.
\]
The cross terms cancel because $M_aM_b=M_cM_d$.
The new expression uses $|A|-2$ squares.
\end{proof}

\paragraph{Prior work on the argument.}
The Gram-matrix method is classical. Choi, Lam, and
Reznick~\cite{clr} (Theorem~2.4 and the discussion before
Example~2.13) identify SOS length with minimum positive semidefinite
Gram rank and discuss uniqueness. Blomenhofer~\cite{blomenhofer}
(Proposition~4.8 and its final sentence) gives a related uniqueness
argument from quadratic independence and forced SOS support.
We make no claim that the method is new. We did not locate the exact
Boolean assignment-polynomial specialization above in the checked
sources.

\begin{lemma}[Binary minimality and pair uniqueness]
\label{lem:sos-minimal-code}
For a binary linear code $C\le\mathbb F_2^N$, pair uniqueness is
equivalent to the following property: if $u,v\in C$, $u\ne0$, and
$\operatorname{supp}(u)\subseteq\operatorname{supp}(v)$, then $u=v$.
\end{lemma}
\begin{proof}
Assume the displayed property and $M_aM_b=M_cM_d$.
If $a=b$, blockwise comparison gives $c=d=a$.
Otherwise put $t=a+b=c+d\ne0$ and $u=a+c\in C$, with addition
in $\mathbb F_2^N$.
At a coordinate where $t$ is zero, the two bits from $a,b$ agree;
the monomial equality forces the bits from $c,d$ to agree with them.
Hence $u$ is zero at that coordinate and
$\operatorname{supp}(u)\subseteq\operatorname{supp}(t)$.
The hypothesis gives $u=0$ or $u=t$.
These alternatives yield $(c,d)=(a,b)$ or $(b,a)$, respectively.

If the displayed property fails, choose $0\ne u\ne v$ with the
specified support inclusion.  The words $0,u,v,u+v$ are four distinct
codewords.  At each coordinate, the unordered pair from $0,v$ equals
the unordered pair from $u,u+v$.  Therefore
$M_0M_v=M_uM_{u+v}$, which violates pair uniqueness.
\end{proof}

\begin{corollary}[Exact code SOS length]
\label{cor:sos-code-exact}
For a binary linear code $C$, the equality
$\operatorname{sos}_{\mathbb R}(P_C)=|C|$ holds if and only if every
nonzero codeword is support-minimal.
\end{corollary}
\begin{proof}
Combine Theorem~\ref{thm:sos-exact-criterion} with
Lemma~\ref{lem:sos-minimal-code}.
\end{proof}

\begin{lemma}[Weight criterion; classical]
\label{lem:sos-code-weights}
If the largest nonzero weight $w_{\max}$ of a binary linear code is
less than twice its minimum nonzero weight $d$, every nonzero
codeword is support-minimal.
\end{lemma}
\begin{proof}
If $0\ne u\ne v$ and
$\operatorname{supp}(u)\subseteq\operatorname{supp}(v)$, then $u$
and $u+v$ are nonzero and have disjoint supports.  Thus
\[
 \operatorname{wt}(v)=\operatorname{wt}(u)+\operatorname{wt}(u+v)
 \ge 2d>w_{\max},
\]
which is impossible.
\end{proof}

This is the binary case of the Ashikhmin--Barg
condition~\cite{ab} (Lemma~2.1(3), printed p.~2010, and Example~2,
printed p.~2011).

\begin{lemma}[Explicit fast minimal codes, from a published code theorem]
\label{lem:sos-fast-minimal}
There is a deterministic polynomial-time-uniform family of injective
linear encoders $E_k:\mathbb F_2^k\to\mathbb F_2^{N_k}$, with
$N_k=\Theta(k)$ and $O(N_k)$ Boolean circuit size, whose image codes
$C_k$ are minimal.  In particular,
\[
 \operatorname{sos}_{\mathbb R}(P_{C_k})=2^k=2^{\Omega(N_k)}.
\]
\end{lemma}
\begin{proof}
Use Brehm, Ishai, and Resch~\cite{bir}, Theorem~19, which restates
Druk and Ishai~\cite{di14}, Theorem~2.
With $q=2$ and $\delta=3/4$, it gives constants $\beta\in\mathbb N$
and $\rho>0$ and codes
$D_k\le(\mathbb F_2^\beta)^{m_k}$ of dimension $k$, rate
$k/(\beta m_k)\ge\rho$, and relative block distance at least $3/4$.
Their encoders have polynomial-time-uniform $O(k)$-wire arithmetic
circuits over $\mathbb F_2$; the circuit and uniformity conventions
are those of that paper's Section~2.2.

Apply to each block the simplex map
\[
 S(z)=\bigl(\langle a,z\rangle\bigr)_{a\in\mathbb F_2^\beta\setminus\{0\}}.
\]
It is injective because the indices $a$ include all coordinate basis
vectors.  If $z\ne0$, choose an index vector $a_0$ with
$\langle a_0,z\rangle=1$.
Translation by $a_0$ pairs the index vectors with dot product zero
with those having dot product one.  Exactly $2^{\beta-1}$ indices
therefore have dot product one, and removing the zero index changes
none of these one entries.

The concatenated code has length $N_k=(2^\beta-1)m_k=\Theta(k)$
and dimension $k$.  Every nonzero word has weight at least
$(3/4)m_k2^{\beta-1}$, and every word has weight at most
$m_k2^{\beta-1}$.  Thus $w_{\max}<2d$ and
Lemma~\ref{lem:sos-code-weights} applies.
The simplex operation has constant size per block because $\beta$
is fixed.  Its composition with the given encoder retains
polynomial-time uniformity and $O(N_k)$ size.
Finally, apply Corollary~\ref{cor:sos-code-exact}.
\end{proof}

\paragraph{Source verification.}
The exact Druk--Ishai input theorem was checked in the published
Brehm--Ishai--Resch restatement, not in the original 2014 full paper.
Theorem~19 is deterministic and must not be confused with that paper's
randomized GV-optimal Theorem~9.  No choice of a good random seed is
used in Lemma~\ref{lem:sos-fast-minimal}.

\begin{theorem}[Signed formulas versus explicit squares, assuming the compiler]
\label{thm:sos-signed-separation}
Assume the integral weighted formula compiler for 3-CNFs of size $s$
with output formula size $2^{O(s/\sqrt{\log s})}$.
There is an explicit family of nonnegative integer polynomials
$P_k$ in $2N_k$ real variables, of total degree $2N_k$ and individual
degree at most two, such that
\begin{enumerate}
\item $P_k$ has a constant-free integer arithmetic formula of size
      $2^{O(N_k/\sqrt{\log N_k})}$;
\item every identity $P_k=\sum_{\ell=1}^{D}p_\ell^2$ with real
      polynomials $p_\ell$ has $D\ge2^{\Omega(N_k)}$.
\end{enumerate}
The formula in the first assertion is syntactically multilinear in
formal squared-weight variables.  After these variables are replaced
by actual squares, only the inserted square gates multiply overlapping
owner blocks.
\end{theorem}
\begin{proof}
Take the encoders and codes from Lemma~\ref{lem:sos-fast-minimal},
and put $P_k=P_{C_k}$.
The relation $z=E_k(u)$ has an $O(N_k)$-size Boolean verification
circuit, counting both $u\in\{0,1\}^k$ and
$z\in\{0,1\}^{N_k}$ among its inputs.
Encode its gate equations and true output by a Tseitin 3-CNF.
There are $O(N_k)$ variables and clauses.  For each message $u$,
there is exactly one satisfying assignment to the gate and codeword
variables: the gates must take their circuit values and $z$ must be
$E_k(u)$.  Injectivity gives exactly one such assignment per codeword.

Assign weights $v_{i,0},v_{i,1}$ to codeword bit $z_i$ and assign both
weights of every message and internal variable the constant one.
The weighted count is the polynomial
\[
 Z_k(v)=\sum_{u\in\{0,1\}^k}\prod_i v_{i,E_k(u)_i}
       =\sum_{c\in C_k}\prod_i v_{i,c_i}.
\]
The assumed compiler supplies a constant-free integer syntactically
multilinear formula of size $2^{O(N_k/\sqrt{\log N_k})}$.
Replace each leaf $v_{i,b}$ by $w_{i,b}w_{i,b}$.
This increases size by a constant factor and computes $P_k$.
Every original multiplication still combines disjoint owner blocks;
only the inserted square gates have overlapping blocks.
The formal degree is at most $2N_k$.
The lower bound is Lemma~\ref{lem:sos-fast-minimal}.

For coefficient explicitness, construct the encoder's generator
matrix by evaluating it on the $k$ basis vectors and use Gaussian
elimination to test whether a requested assignment lies in its image.
These computations take polynomial time in $N_k$.
The coefficient is one exactly for the squared assignment monomials
of codewords, and is zero otherwise.
\end{proof}

\begin{corollary}[No polynomial-size universal SOS conversion, assuming the compiler]
\label{cor:sos-no-conversion}
There is no polynomial-overhead conversion that turns every
constant-free signed formula obtained above into an explicitly listed
sum of real polynomial squares computing the same polynomial.
The obstruction persists when the computations of the square roots
may share subcircuits.
\end{corollary}
\begin{proof}
A fixed polynomial in the upper size bound is still
$\exp(O(N_k/\sqrt{\log N_k}))$.
The output must explicitly contain at least
$\exp(\Omega(N_k))$ square summands by
Theorem~\ref{thm:sos-signed-separation}.
Sharing the computations of the square roots does not alter their
number, which is the quantity bounded by the theorem.
\end{proof}

\subsection{Short EF proofs of bounded-integer positivity}

The preceding lower bound does not obstruct the direct EF argument.
The following lemma extends negative descent to nonnegative weights;
its proof uses no explicit sum-of-squares representation.

\begin{lemma}[Weighted negative descent, from the EF arithmetic interface]
\label{lem:sos-weighted-descent}
Let $W$ be a smooth set-multilinear integer formula of size $S$ on
$N$ weight pairs.  Let $f$ be a Boolean-valued circuit, and suppose
an EF proof of length $T$ establishes
$W(\operatorname{pin}(b))=f(b)$ for arbitrary Boolean bits $b$.
For arbitrary unsigned $B$-bit weights $w_{i,0},w_{i,1}$ and
arbitrary Boolean bits $a$, there are EF proofs of
\[
 W(w)\ge f(a)M_a(w),
 \qquad\text{and therefore}\qquad W(w)\ge0,
 \tag{$\dagger$}\label{eq:sos-weighted-inequality}
\]
of size polynomial in $S,N,B,T$, and the circuit size of $f$.
The constants in $W$ are assumed constant-free, or supplied with
bit lengths included in these parameters.
\end{lemma}
\begin{proof}
Smooth set-multilinearity gives the weighted splitting identity
\[
 Q(w)=w_{i,0}Q|_{(w_{i,0},w_{i,1})=(1,0)}
      +w_{i,1}Q|_{(w_{i,0},w_{i,1})=(0,1)}.
 \tag{$\ddagger$}\label{eq:sos-weighted-splitting}
\]
It has a polynomial-size arithmetic derivation by induction on the
formula.  At a weight leaf it is $w_{i,0}=w_{i,0}\cdot1+w_{i,1}\cdot0$
or the counterpart for $w_{i,1}$; at a zero leaf both sides are zero.
At a sum, add the two child identities and reassociate.
At a product, only one child contains pair $i$; apply its identity
and distribute over the other child.  The other child is unchanged.

Assume the first inequality in~\eqref{eq:sos-weighted-inequality}
fails, and put
\[
 Q_a(w)=W(w)-f(a)M_a(w),\qquad
 M_a(w)=\prod_i\bigl((1-a_i)w_{i,0}+a_iw_{i,1}\bigr).
\]
This is smooth set-multilinear in the weight pairs, with $a$ as
parameters, and $Q_a(w)<0$.
Suppose a prefix $b_1,\ldots,b_{j-1}$ has been selected so that
the partial value remains negative, while all remaining pairs
retain their originally supplied weights.  Let $u_j,v_j$ be the
partial values obtained by additionally pinning pair $j$ to zero
or one.  Equation~\eqref{eq:sos-weighted-splitting} says that the
current negative value is $w_{j,0}u_j+w_{j,1}v_j$.

If $u_j<0$, select $b_j=0$.
If $u_j\ge0$, nonnegativity of $w_{j,0}$ gives
$w_{j,0}u_j\ge0$, hence $w_{j,1}v_j<0$.
Nonnegativity of $w_{j,1}$ then forces $v_j<0$, so select $b_j=1$.
The second case cannot occur with $w_{j,1}=0$ because it would
give $0<0$.  Thus the argument includes zero weights.

After all pairs have been pinned, the presumed negative value is
\[
 Q_a(\operatorname{pin}(b))=f(b)-f(a)[a=b].
\]
When $a=b$, this is zero by gate-by-gate equality propagation
through $f$.  When $a\ne b$, it is the Boolean value $f(b)$.
Both alternatives contradict negativity.

Implement each selection by extension variables for the two
evaluations and the comparison $u_j<0$.
These definitions are acyclic: step $j$ refers only to $a$, the
original weights, and already defined prefix bits.
There are $O(N)$ evaluations of formulas of size polynomial in
$S+N$, and $O(N)$ instances of the weighted splitting proof.
All input values have $O(B+1)$ bits; the formula inequalities
$\log_2(|x+y|+1)\le1+\max\{\log_2(|x|+1),\log_2(|y|+1)\}$ and
$\log_2(|xy|+1)\le\log_2(|x|+1)+\log_2(|y|+1)$ bound each
intermediate word length by a polynomial in $(S+N)(B+1)$.
The EF arithmetic interface proves the displayed ring and order
steps with polynomial overhead in this width.
This gives the first inequality.
Every factor in $f(a)M_a(w)$ is nonnegative, proving the second.
\end{proof}

\begin{corollary}[Explicit squares versus EF positivity, from the pinning identity]
\label{cor:sos-ef-positivity}
For the family $P_k$ in
Theorem~\ref{thm:sos-signed-separation}, the small signed formulas
can be chosen so that their nonnegativity on arbitrary signed
$B$-bit integer inputs has EF proofs of size
\[
 2^{O(N_k/\sqrt{\log N_k})}(B+1)^{O(1)}.
\]
For the same polynomials, every real polynomial sum-of-squares
identity still needs $2^{\Omega(N_k)}$ square summands.
This compares bounded-integer EF positivity proofs with explicit
real polynomial SOS identities; it is not a lower bound for every
proof system called sum-of-squares.
\end{corollary}
\begin{proof}
Use the Boolean network for the encoder-verification relation in
Theorem~\ref{thm:sos-signed-separation}, and apply the all-port
weighted compiler after subdivision.
For each designated codeword bit choose one port carrying that
bit, and assign its weight pair the desired codeword weights.
Assign both weights at every other port the constant one.
Every consistent network assignment uses each chosen codeword
weight once.  The message and internal gate values are uniquely
determined, so the specialized weighted polynomial is $Z_k$ from
that theorem.
The network size and number of ports remain $O(N_k)$.

The all-port pinning theorem and its quantitative derivation give
an EF proof of the pinning identity of length
$2^{O(N_k/\sqrt{\log N_k})}$.
Apply Lemma~\ref{lem:sos-weighted-descent} before making the
specialization of the port weights.
Now replace each chosen codeword weight by the square of the
corresponding signed integer input and all other weights by one.
To prove that such a square is nonnegative, split on the sign of
its input $z$.  If $z\ge0$, multiplication preserves nonnegativity,
so $z^2\ge0$.  If $z<0$, the order rules give $-z>0$, the first
case gives $(-z)^2\ge0$, and the ring identities give
$(-z)^2=z^2$.  These are instances of the imported EF arithmetic
library.  A full-result multiplication uses $O(B+1)$ bits, so each
substituted weight has that word length.
The resulting polynomial is $P_k$ and the asserted proof-size
bound follows.  Corollary~\ref{cor:sos-code-exact} supplies the
independent SOS-length lower bound.
\end{proof}

\paragraph{Hand check on a four-clause formula.}
Let $C=\{000,011,101,110\}$.
All its nonzero words have weight two, so it is minimal.
Its four assignment monomials are
\[
 a_0b_0c_0,\quad a_0b_1c_1,\quad a_1b_0c_1,\quad a_1b_1c_0.
\]
The six mixed products are
\[
\begin{array}{lll}
a_0^2b_0b_1c_0c_1,&a_0a_1b_0^2c_0c_1,&a_0a_1b_0b_1c_0^2,\\
a_0a_1b_0b_1c_1^2,&a_0a_1b_1^2c_0c_1,&a_1^2b_0b_1c_0c_1.
\end{array}
\]
They are distinct, and none is a diagonal product because each has
an odd exponent.  Evaluation at the four assignment points outside
$C$ kills the other four possible assignment monomials in every
square root.  Thus every Gram matrix is $I_4$, and four squares are
necessary and sufficient.
The code is precisely the satisfying set of
\[
 (a\vee b\vee\neg c)\wedge(a\vee\neg b\vee c)
 \wedge(\neg a\vee b\vee c)\wedge(\neg a\vee\neg b\vee\neg c).
\]
These clauses exclude $001,010,100,111$, respectively.
The encoder $(s,t)\mapsto(s,t,s\oplus t)$ gives the same four words,
so its weighted verification relation gives the claimed polynomial.

\paragraph{Hand check when minimality fails.}
For $C=\mathbb F_2^2$, the collision
$M_{00}M_{11}=M_{01}M_{10}$ yields
\[
 (a_0^2+a_1^2)(b_0^2+b_1^2)
 =(a_0b_0+a_1b_1)^2+(a_0b_1-a_1b_0)^2.
\]
Thus four assignment squares compress to two.  The signs inside the
square roots are precisely what makes the compression possible here.

\paragraph{Limits of the conclusion.}
The theorem blocks a universal replacement of the counting
polynomial by a short explicit squared norm, even allowing signed
polynomials as its coordinates.  It does not block Frege reasoning
about positivity, implicit sums, local identities, rational functions,
Boolean quotient rings, or a fixed numerical specialization.
In particular, an unsatisfiable formula has zero weighted polynomial,
which needs no nonzero square summands.

The theorem also does not trigger the arithmetic circuit lower-bound
implications of Hrube\v{s}--Wigderson--Yehudayoff~\cite{hwy} or
Rossman--Zhu~\cite{rz}. Theorem~1.3 of~\cite{rz} requires
sum-of-squares lower bounds over algebraically closed or finite fields
of characteristic different from two.  Its consequence
$\mathsf{VNC}^1\ne\mathsf{VNP}$ requires $d$ blocks of $n$ variables
with $\omega(1)\le d\le O(\log n/\log\log n)$.
Here the blocks have $n=2$ and $d=N_k$, and the proof uses the order
property of $\mathbb R$.  The crucial implication
$\sum p_i^2=0\Rightarrow p_i=0$ fails over $\mathbb C$, as witnessed
by $1^2+\mathrm{i}^2=0$.
Their Section~1.1 explicitly records that the analogous implications
are not known over $\mathbb R$.


\section{A coefficient that is hard in every variable order}
\label{sec:all-orders-coefficient-barrier}

\paragraph{Status.}
The statements below are proved, using Theorem~9 of Brehm, Ishai,
and Resch~\cite{bir} as a coding-theoretic input. They concern
coefficient-by-coefficient normalization, defined below, into the
sorted form for one variable order. Normalization into an order
prescribed in advance is the special case of an order that does not
depend on the input; the statements below show that choosing the
normalization order adaptively does not repair coefficient-by-coefficient
normalization. They do not give a lower bound for Frege or for
noncommutative IPS.
The unsatisfiable formulas in the construction contain the two unit
clauses $u$ and $\neg u$. Novelty relative to the circuit and branching
program literature has not been established.

Lemma~\ref{lem:code-rank-unique-witnesses} gives a rank lower bound
for code indicators, Lemma~\ref{lem:evaluation-rank-free-formula}
converts rank into a noncommutative formula lower bound, and
Theorem~\ref{thm:uniform-all-orders-coefficient} embeds the result in
an explicit unsatisfiable CNF. The first-failed-clause coefficient
used there is displayed in its proof.

\paragraph{Terminology.}
The Ideal Proof System (IPS) of Grochow and Pitassi~\cite{gp} refutes
a system of polynomial equations $F_1=\cdots=F_r=0$, which includes
the Boolean axioms $X_a^2-X_a$, by an \emph{IPS certificate}: a
polynomial $C(X,Y_1,\ldots,Y_r)$ with $C(X,0,\ldots,0)=0$ and
$C(X,F_1,\ldots,F_r)=1$, measured by its circuit or formula size. In
\emph{noncommutative IPS} the certificate is a noncommutative (free)
formula, and the commutators $X_aX_b-X_bX_a$ are added as axioms.
The commutator and Boolean-axiom placeholders in
Corollary~\ref{cor:adaptive-normalization-barrier} stand for these
axioms.
For a CNF whose clauses are listed in a fixed order, every falsifying
assignment has a first falsified clause. The
\emph{first-failed-clause coefficient} of a clause $\sigma$ is the
multilinear indicator, over $\mathbb F_2$ and in the variables that do
not occur in $\sigma$, of the assignments to those variables that
satisfy every clause listed before $\sigma$ when the variables of
$\sigma$ take the unique values falsifying $\sigma$. The
\emph{contribution} $\Gamma_\sigma$ of $\sigma$ is this coefficient
multiplied by the indicator that $\sigma$ is falsified.
\emph{Coefficient-by-coefficient normalization} starts from a free
formula for each coefficient and rewrites it separately into
$\operatorname{sort}_\pi$ of that coefficient, in the sense defined
next, for some order $\pi$ that may differ between coefficients,
before the contributions are added.

All polynomials and matrix ranks in this section are over $\mathbb
F_2$. For a set $B\subseteq\{0,1\}^n$, put
\[
 I_B(X)=\sum_{b\in B}\prod_{j=1}^n\eta_{X_j}(b_j),
 \qquad \eta_X(0)=1+X,\quad \eta_X(1)=X.
\]
This is a multilinear polynomial. On a Boolean input $a$, every
summand except the one indexed by $a$ evaluates to zero, and that
summand, when present, evaluates to one. Thus $I_B$ is the Boolean
indicator of $B$. For an order $\pi$, let
$\operatorname{sort}_{\pi}(I_B)$ be its canonical free lift: replace
each squarefree commutative monomial by the word containing those
variables in order $\pi$.

\begin{lemma}[Rank survives arbitrary unique witnesses]
\label{lem:code-rank-unique-witnesses}
Let $C\le\mathbb F_2^N$ be a binary linear code of minimum distance
$d$ whose dual has minimum distance $d^\perp$. Let
$g:C\longrightarrow\{0,1\}^w$ be any function, with no linearity
assumption, and let
\[
 B=\{(c,g(c)):c\in C\}\subseteq\{0,1\}^{n},\qquad n=N+w.
\]
Fix an integer $1\le k<\min\{d,d^\perp\}$ and any order $\pi$ of
the $n$ coordinates. End a prefix $P$ of $\pi$ immediately after
the $k$th code coordinate; the prefix may also contain witness
coordinates. Put $Q=[n]\setminus P$. The evaluation matrix
\[
 E_{a,b}=I_B(a,b),\qquad
 a\in\{0,1\}^{P},\ b\in\{0,1\}^{Q},
\]
has rank at least $2^k$.
\end{lemma}
\begin{proof}
Write $J=P\cap[N]$, so $|J|=k$. The projection of $C$ onto $J$
is all of $\mathbb F_2^J$. Indeed, if its image were a proper
subspace, row reduction would produce a nonzero vector of
$\mathbb F_2^J$ orthogonal to that image: a homogeneous linear
system of rank less than $k$ has a nonzero solution. Extending
that vector by zero outside $J$ would give a nonzero word of
$C^\perp$ supported on at most $k$ coordinates. This contradicts
$k<d^\perp$.

There is therefore an extendible prefix assignment for each of the
$2^k$ assignments to $J$. Prefix assignments obtained from different
assignments to $J$ are distinct. Hence $E$ has at least $2^k$
nonzero rows.

Two distinct nonzero rows cannot have a common nonzero column.
Suppose otherwise that $(a,b)$ and $(a',b)$ both belong to $B$.
Let their codewords be $c,c'$. They agree on every code coordinate
outside $J$, so $c-c'$ has weight at most $k$. Since $k<d$, this
forces $c=c'$. The witness is a function of the codeword, so
$g(c)=g(c')$, giving $(a,b)=(a',b)$ and hence $a=a'$, a
contradiction.

Select any $2^k$ nonzero rows. Their supports are nonempty and
pairwise disjoint. In any linear combination of these rows, a
column in the support of a selected row receives only that row's
coefficient. A zero linear combination therefore has every
coefficient zero. These rows are linearly independent.
\end{proof}

\begin{lemma}[From evaluation rank to a free-formula lower bound]
\label{lem:evaluation-rank-free-formula}
Let $f$ be a multilinear polynomial in $n$ variables. Let $P,Q$
partition its variables, with all variables in $P$ before those in
$Q$ in an order $\pi$. If the Boolean evaluation matrix of $f$
across this partition has rank $R$, then every noncommutative
formula for $\operatorname{sort}_{\pi}(f)$ has size at least
\[
 \frac{R}{(|P|+1)(|Q|+1)}\ge\frac{R}{(n+1)^2}.
\]
The formula is not required to be homogeneous or syntactically
ordered, and cancellations inside it are allowed.
\end{lemma}
\begin{proof}
Write
\[
 f=\sum_{A\subseteq P,\ B\subseteq Q}c_{A,B}X_AX_B
\]
and let $K=(c_{A,B})$. The Boolean evaluation matrix satisfies
\[
 E=Z_PKZ_Q^{\mathsf T},\qquad
 (Z_P)_{a,A}=\prod_{i\in A}a_i.
\]
Identify a Boolean assignment $a$ with its support. Then
$(Z_P)_{a,A}=1$ exactly when $A\subseteq a$. Ordering subsets by
cardinality, with the same tie-breaking on rows and columns,
makes this matrix lower triangular with diagonal entries one.
Thus $Z_P$ is invertible; the same argument applies to $Z_Q$.
Consequently $\operatorname{rank}K=\operatorname{rank}E=R$.

Suppose a free formula of size $S\ge1$ computes
$\operatorname{sort}_{\pi}(f)$. Convert it into a source--target
acyclic algebraic branching program as follows. A leaf is one
edge bearing its scalar or variable. An addition takes disjoint
child programs and joins new source and target vertices by edges
labelled one. A multiplication joins the left target to the right
source by an edge labelled one. The resulting program has at most
$S$ variable-labelled edges, and its path polynomial is the formula
polynomial, as follows inductively from the two constructions.

For integers $a\ge0,b\ge1$, form the coefficient matrix whose rows
are words of length $a$, columns are words of length $b$, and whose
entry at $(v,w)$ is the coefficient of $vw$. Every contributing
source--target path of degree $a+b$ has a unique $(a+1)$st variable
edge. For each possible distinguished edge, the sum of path
contributions factors as a column vector of prefix coefficients
times a row vector of suffix coefficients, with the distinguished
edge included in the suffix. This is an outer product. Concatenating
any such prefix and suffix is valid: if they shared another vertex,
their concatenation would yield a directed cycle in the acyclic
program. Summing over at most $S$ variable edges proves that this
coefficient matrix has rank at most $S$.

For fixed $a,b$, restrict this matrix to the increasing words for
sets $A\subseteq P$ of size $a$ and $B\subseteq Q$ of size $b$.
It is the corresponding block of $K$, so that block has rank at
most $S$. When $b=0$, the block has one column and also has rank
at most $1\le S$. Embed each of these blocks in the full row and
column index sets, with zeros outside the block. Their sum is
$K$. There are $(|P|+1)(|Q|+1)$ blocks, and the rank of a sum is
at most the sum of the ranks, because its column space lies in the
sum of the summands' column spaces. Therefore
\[
 R\le (|P|+1)(|Q|+1)S,
\]
which is the claimed bound.
\end{proof}

\begin{theorem}[Every order is expensive, even with unique witnesses]
\label{thm:all-orders-code-indicator}
In the setting of Lemma~\ref{lem:code-rank-unique-witnesses}, every
order $\pi$ satisfies
\[
 \operatorname{FormulaSize}_{\rm nc}
       (\operatorname{sort}_{\pi}(I_B))
 \ \ge\ \frac{2^k}{(n+1)^2}
 \qquad (1\le k<\min\{d,d^\perp\}).
\]
In particular, if $d,d^\perp\ge\delta N$ for a fixed
$\delta>0$ and $w=O(N)$, the bound is $2^{\Omega(n)}$,
simultaneously for every variable order.
\end{theorem}
\begin{proof}
For the chosen order, take the prefix ending at the $k$th code
coordinate. The first lemma gives evaluation rank at least $2^k$
for this prefix. The second lemma gives the displayed formula-size
bound. In the final assertion choose
$k=\lfloor\delta N/2\rfloor$ for all sufficiently large $N$.
Then $k<d,d^\perp$, and $n=O(N)$ makes the denominator polynomial
and $k=\Omega(n)$.
\end{proof}

\begin{theorem}[A uniform sparse contradiction with an all-orders-hard coefficient]
\label{thm:uniform-all-orders-coefficient}
There is a polynomial-time constructible sequence of unsatisfiable
CNFs $U_i$ of width at most three, with $n_i$ variables and $O(n_i)$
clauses, and a specified first-failed-clause coefficient $Q_i$,
such that every noncommutative formula for
$\operatorname{sort}_{\pi}(Q_i)$ has size $2^{\Omega(n_i)}$ for
every variable order $\pi$. Each $U_i$ contains the two contradictory
unit clauses $u$ and $\neg u$ and consequently has an immediate
Frege refutation.
\end{theorem}
\begin{proof}
Use~\cite[Theorem~9]{bir} with $q=2$, $R=1/2$,
and $\epsilon=1/8$. Its computational model, defined in
\S2.2, uses binary arithmetic gates, measures size in wires,
and defines uniformity by polynomial-time generation of the
circuits. The theorem supplies a uniform randomized systematic
encoder
\[
 E_i(r,m)=(m,P_i(r,m))\in\mathbb F_2^{N_i}
\]
of size $O(N_i)$, with $O(N_i)$ random input bits $r$.
For each sufficiently large $i$ there exists a seed $r_i^*$ such
that the resulting linear code $C_i$ and its dual both have
minimum distance at least $\delta N_i$ for a fixed $\delta>0$.
Indeed, the theorem bounds each failure probability by
$2^{-\epsilon N_i}$, uses the same randomness for the primal
and dual codes, and gives their relative distances bounded below
by the positive quantities
$h_2^{-1}(1-R_i-\epsilon)$ and
$h_2^{-1}(R_i-\epsilon)$, with $R_i\to1/2$.
For large $i$, both are at least $h_2^{-1}(1/4)>0$ and
$2\cdot2^{-\epsilon N_i}<1$.

Treat $r$ as ordinary input variables, rather than choosing a seed
in the construction. Over $\mathbb F_2$, addition and multiplication
are XOR and AND, and each has a fixed-size Boolean implementation.
Construct the width-three CNF $G_i(r,m,y,z)$ that expresses the
gate equations of the encoder and the equalities
$y=P_i(r,m)$. Here $z$ contains one variable for each internal
gate. For an AND gate $z=a\wedge b$ use
\[
 (\neg z\vee a),\quad(\neg z\vee b),\quad
 (z\vee\neg a\vee\neg b).
\]
For an XOR gate $z=a\oplus b$ use
\[
 (a\vee b\vee\neg z),\quad
 (a\vee\neg b\vee z),\quad
 (\neg a\vee b\vee z),\quad
 (\neg a\vee\neg b\vee\neg z).
\]
Equality of an output with $y_j$ uses the two implication clauses.
Fixed scalar wires are replaced by their bits. Each set of clauses
forces exactly the named gate function, by checking its four input
pairs. Processing the encoder in topological order proves that
every input $(r,m)$ has exactly one gate assignment $z$, and then
exactly one compatible output $y$. The number of variables and
clauses is $O(N_i)$, while the code coordinates $(m,y)$ themselves
already number $N_i$. All descriptions are generated in polynomial
time from the encoder description.

Define $U_i$ by listing all clauses of $G_i$, then $u$, then
$\neg u$. Its coefficient for first failure at the clause $u$ is
\[
 Q_i(r,m,y,z)=
 \sum_{a\models G_i}\prod_{v\in\{r,m,y,z\}}\eta_v(a_v)
 =I_{\operatorname{SAT}(G_i)}.
\]
The omitted variable $u$ is forced to zero by failure of that unit
clause. Thus it does not enter the coefficient.

Now fix the seed variables to $r_i^*$ only in the lower-bound
argument. The restricted polynomial is the Boolean indicator of
\[
 B_i=\{(c,g_i(c)):c\in C_i\},
\]
where $c=(m,y)$ and $g_i(c)$ is the unique gate assignment.
There are $O(N_i)$ witness coordinates. To justify that the
restriction is exactly this indicator as a polynomial, first note
that it is multilinear and has the claimed Boolean evaluations.
The triangular matrix argument in
Lemma~\ref{lem:evaluation-rank-free-formula}, with all variables
on one side, says a multilinear polynomial is uniquely determined
by those evaluations.

Scalar restriction commutes with canonical free lifting: deleting
variables assigned one and deleting monomials containing variables
assigned zero leaves each surviving word in the induced variable
order, and the summed coefficients equal the scalar specialization
of the commutative polynomial. Therefore any free formula for
$\operatorname{sort}_{\pi}(Q_i)$ restricts, without increasing
size, to a formula for
$\operatorname{sort}_{\pi'}(I_{B_i})$, where $\pi'$ is the
induced order after deleting the seed variables. The previous
theorem gives size $2^{\Omega(N_i)}=2^{\Omega(n_i)}$.

Finally, the unit clauses $u$ and $\neg u$ derive a contradiction
by one application of the fixed propositional contradiction rule
or a fixed Frege derivation. This establishes the last assertion
without using any property of $G_i$.
\end{proof}

\begin{corollary}[Adaptive coefficient normalization is not a cheap repair]
\label{cor:adaptive-normalization-barrier}
For the coefficients $Q_i$ above, let $A_i$ be any free formula
whose commutative image is $Q_i$ and whose size is
$2^{o(n_i)}$. Select an arbitrary order $\pi_i$, even depending
on the entire input and on $A_i$. Suppose a free formula $R_i$
with commutator placeholders $Z_{ab}$ and Boolean-axiom
placeholders $B_a$ satisfies the exact free-algebra identity
\[
 A_i+R_i(X,(X_aX_b-X_bX_a)_{a,b},(X_a^2-X_a)_a)
   =\operatorname{sort}_{\pi_i}(Q_i).
\]
Then $R_i$ has size $2^{\Omega(n_i)}$.
\end{corollary}
\begin{proof}
Every placeholder is replaced by a formula of constant size.
As a formula is a tree, replacing all its placeholder leaves
increases size by at most a fixed multiplicative constant.
The left-hand side therefore has formula size
$O(|A_i|+|R_i|)$. The preceding theorem lower-bounds this by
$2^{cn_i}$ for some $c>0$ and all sufficiently large $i$.
Because $|A_i|=2^{o(n_i)}$, it is at most $2^{cn_i}/2$ for all
sufficiently large $i$, after absorbing fixed constants. The
remaining contribution forces $|R_i|=2^{\Omega(n_i)}$.
\end{proof}

\paragraph{Hand check with a witness preceding a code coordinate.}
Take $C=\{000,111\}$, which has minimum distance $3$ and dual
minimum distance $2$. Let $g(c)=c_1c_2$, so
$B=\{0000,1111\}$, with witness coordinate $z$. Use order
$z,x_2,x_1,x_3$ and $k=1$. The prescribed prefix is
$P=\{z,x_2\}$. Restricted to rows $00,11$ and columns $00,11$,
the evaluation matrix is
\[
 \begin{pmatrix}1&0\\0&1\end{pmatrix}.
\]
The two rows are disjoint because the two codewords cannot agree
on the two public coordinates outside the prefix. In polynomial
form the indicator is
\[
 (1+x_1)(1+x_2)(1+x_3)(1+z)+x_1x_2x_3z.
\]
The degree-four terms cancel in characteristic two. Thus the
argument genuinely applies to a polynomial with cancellation and
with an auxiliary witness placed before a public coordinate.

\paragraph{Exact implication for the global residual.}
This result rules out converting each first-failed coefficient to
one order before adding the clause contributions, even if the
orders are chosen separately and adaptively. It does not rule out
correcting the global residual
$H=\sum_\sigma\Gamma_\sigma-1$, the sum of all clause contributions
minus one. In the displayed CNFs, a global
certificate can discard the entire $G_i$ prefix and use the
contradictory unit clauses. A valid construction for the global
residual may therefore exploit cancellation between coefficients
or choose a different IPS certificate. The theorem gives no lower
bound against either option.



\section{Verification record}
\label{sec:verification}

\paragraph{Proof checks.}
The pinning, merge reconstruction, and size accounting were
checked independently in this research pass. The checks included
the placement of port weights after subdivision, the complete
split sum before orientation removal, the packing bijection,
the bit-width bound via expanded formulas, and the acyclic
negative-descent definitions. The sections include hand checks
on an edge, a one-bit crossing, a $2\times2$ merge, two selector
branches, and four-clause formulas. These checks supplement
the proofs; they are not substitutes for them. The new proof
has not been machine-verified.

\paragraph{Primary sources.}
The references below give the versions and locators actually used.
The supplied manuscript and Szeider version 2 were read locally.
The Krapivin--Li--Shi manuscript was downloaded from its STOC
submission link after the ePrint body was inaccessible. The
original Druk--Ishai paper was not obtained; its exact coding
theorem was checked in the stated Brehm--Ishai--Resch restatement.

\paragraph{Acknowledgments.}
AI was used to produce the results, proofs, and text of this note.

\begin{thebibliography}{99}
\bibitem{cruz-ntime} Charlie Cruz. \emph{Separating Nondeterministic Time from
Deterministic Space}. Supplied
\texttt{separating-ntime-from-space.tex}.
Locators: \texttt{lem:codes}, \texttt{lem:top},
\texttt{lem:atone}, \texttt{cor:cuts}, \texttt{lem:transport},
\texttt{sec:merge}, \texttt{sec:formula}, and \texttt{lem:paths}.

\bibitem{szeider} Stefan Szeider. Supplied note on consequences of the ETH
refutations, version 2 of November 2026. \S7, Open Problems
1 and 3, printed p.~10; paragraph on the soundness invariant,
p.~11. The public announcement still states the EF consequence
with the IPS simulation hypothesis.
\url{https://ethdiscussion.blog/2026/11/01/announce-your-results/}
and the supplied Zenodo identifier
\url{https://doi.org/10.5281/zenodo.24299623}.

\bibitem{kls} Andrew Krapivin, Jiatu Li, Elaine Shi.
\emph{Fine-Grained Proof Complexity for Elementary Arithmetic in
Extended Frege}. STOC 2027 submission 1284; ePrint 2026/2758.
Appendix A.1, pp.~43--44; Lemma C.1, pp.~63--65, equations
(62)--(76); Theorem A.2, p.~54; signed extension in Appendix J.3,
pp.~138--140. Submission status is not acceptance.
\url{https://eprint.iacr.org/2026/2758}
and
\url{https://github.com/sigact-stoc/stoc2027/releases/download/papers-2/1284.pdf}.

\bibitem{aadk} Noel Arteche, Albert Atserias, Susanna F. de Rezende,
Erfan Khaniki. \emph{The Proof Analysis Problem}, June 2025
version. \S2.3.4, Theorem 2.3, printed p.~18 (PDF p.~20),
states Cook's correspondence theorem; it points to
Kraj\'i\v{c}ek's 2019 book, Theorem 12.4.2, for the proof.
\url{https://narteche.github.io/files/papers/PAP_june2025.pdf}.

\bibitem{BSW} Eli Ben-Sasson and Avi Wigderson.
\emph{Short proofs are narrow---resolution made simple}. J. ACM, 2001.
\href{https://doi.org/10.1145/375827.375835}{Article}.

\bibitem{BI} Eli Ben-Sasson and Russell Impagliazzo.
\emph{Random CNFs are hard for the polynomial calculus}.
\href{https://www.ias.edu/sites/default/files/math/csdm/10-11/ebensasson_random_cnfs_are_hard.pdf}{Paper}, Theorem~1.2.

\bibitem{AH} Albert Atserias and Tuomas Hakoniemi.
\emph{Size-degree trade-offs for sums-of-squares and Positivstellensatz proofs}.
CCC, 2019. \href{https://doi.org/10.4230/LIPIcs.CCC.2019.24}{Proceedings}.

\bibitem{calabro} Chris Calabro. \emph{The Exponential Complexity of
Satisfiability Problems}. UC San Diego, 2009. Chapter~5, equation~(5.4),
Lemmas~5.13--5.14 and Theorem~5.15.
\href{https://cseweb.ucsd.edu/~ccalabro/thesis.pdf}{Dissertation}.

\bibitem{ipz} Russell Impagliazzo, Ramamohan Paturi, Francis Zane.
\emph{Which Problems Have Strongly Exponential Complexity?}
Journal of Computer and System Sciences 63 (2001), 512--530.
\S2, Theorem 1 and \textsc{Reduce}, pp.~520--521;
Lemma 7 and Corollary 1, p.~523.
\url{https://cseweb.ucsd.edu/~paturi/myPapers/pubs/ImpagliazzoPaturiZane_2001_jcss.pdf}.

\bibitem{ccs} Sravanthi Chede, Leroy Chew, Anil Shukla.
\emph{Circuits, Proofs and Propositional Model Counting}.
ECCC TR24-081, revision 1, August 2024; FSTTCS 2024.
Definition 19, Lemma 21 and Theorem 23, PDF pp.~14--15
of the ECCC revision.
\url{https://eccc.weizmann.ac.il/report/2024/081/revision/1/download/}.

\bibitem{krajicek} Jan Kraj\'i\v{c}ek. \emph{Propositional proof
complexity I}, lecture-note draft. \S3.3, Definition 3.3.1,
Lemma 3.3.2, Theorem 3.3.3 and its following discussion,
printed pp.~52--54, for extension variables and circuit-style
proofs.
\url{https://msekce.karlin.mff.cuni.cz/~krajicek/ds1.pdf}.

\bibitem{bir} Martijn Brehm, Yuval Ishai, Nicolas Resch.
\emph{Fast Bounded-Independence Functions and Their Duals}.
ITC 2026, DOI 10.4230/LIPIcs.ITC.2026.8.
\S2.2 specifies uniformity and wire count.
Theorem 9 supplies the randomized systematic code used in
the coefficient appendix. Theorem 19 restates Druk--Ishai
(ITCS 2014), Theorem 2, and supplies the deterministic
constant-block encoder used in the SOS appendix.
\url{https://drops.dagstuhl.de/storage/00lipics/lipics-vol385-itc2026/html/LIPIcs.ITC.2026.8/LIPIcs.ITC.2026.8.html}.

\bibitem{ab} Alexei Ashikhmin, Alexander Barg.
\emph{Minimal Vectors in Linear Codes}. IEEE Transactions on
Information Theory 44(5) (1998), 2010--2017.
Lemma 2.1(3),(5), p.~2010; Example 2, p.~2011.
\url{https://user.eng.umd.edu/~abarg/reprints/MinimalVectors98.pdf}.

\bibitem{clr} Man-Duen Choi, T. Y. Lam, Bruce Reznick.
\emph{Sums of Squares of Real Polynomials}. Proceedings of
Symposia in Pure Mathematics 58.2 (1995).
Theorem 2.4 and the discussion before Example 2.13,
printed pp.~108--110, for the classical Gram-matrix viewpoint.
\url{https://www.math.ucdavis.edu/~deloera/MISC/LA-BIBLIO/trunk/ReznickBruce/choilamreznicksosrealgrammatrix.pdf}.

\bibitem{blomenhofer} Alexander Taveira Blomenhofer.
\emph{On Uniqueness of Power Sum Decomposition}.
SIAM Journal on Applied Algebra and Geometry 9(1) (2025),
211--234. Proposition 4.8 and its concluding observation,
p.~225, give related Gram uniqueness criteria.
\url{https://ir.cwi.nl/pub/35868/35868.pdf}.

\bibitem{rz} Benjamin Rossman, Davidson Zhu.
\emph{Multiquadratic Sum-of-Squares Lower Bounds Imply
VNC$^1\ne$VNP}. arXiv:2512.01227, consulted HTML version.
\S1.1, the qualification following Theorem 1.2;
\S1.2, Theorem 1.3; \S2.1, field conventions.
\url{https://arxiv.org/html/2512.01227v1}.

\bibitem{cook75} Stephen A. Cook.
\emph{Feasibly Constructive Proofs and the Propositional Calculus}.
Proceedings of the 7th Annual ACM Symposium on Theory of Computing
(STOC 1975), 83--97.

\bibitem{di14} Erez Druk, Yuval Ishai.
\emph{Linear-Time Encodable Codes Meeting the Gilbert--Varshamov
Bound and Their Cryptographic Applications}. Proceedings of the 5th
Innovations in Theoretical Computer Science Conference (ITCS 2014),
169--182. Not obtained; see the note on primary sources above.

\bibitem{hwy} Pavel Hrube\v{s}, Avi Wigderson, Amir Yehudayoff.
\emph{Non-Commutative Circuits and the Sum-of-Squares Problem}.
Journal of the American Mathematical Society 24(3) (2011), 871--898.

\bibitem{gp} Joshua A. Grochow, Toniann Pitassi.
\emph{Circuit Complexity, Proof Complexity, and Polynomial Identity
Testing: The Ideal Proof System}. Journal of the ACM 65(6) (2018),
Article 37.
\end{thebibliography}

\end{document}

